\documentclass[11pt,a4paper]{article}
\usepackage[utf8]{inputenc}
\usepackage[T1]{fontenc}
\usepackage[spanish,es-nodecimaldot,es-noquoting]{babel}
\usepackage[scaled=0.95]{helvet}
\renewcommand{\familydefault}{\sfdefault}
\usepackage{amsmath,array,tabularx,booktabs,enumitem,graphicx}
\usepackage[margin=20mm,top=21mm,bottom=21mm,headheight=15pt]{geometry}
\usepackage{xcolor,tcolorbox,fancyhdr,lastpage,microtype,etoolbox,xurl}
\usepackage[hidelinks,unicode]{hyperref}
\ifdefined\pdfinfoomitdate\pdfinfoomitdate=1\fi
\ifdefined\pdftrailerid\pdftrailerid{}\fi
\ifdefined\pdfsuppressptexinfo\pdfsuppressptexinfo=15\fi
\tcbuselibrary{breakable,skins}
\BeforeBeginEnvironment{tabularx}{\par\noindent}
\AfterEndEnvironment{tabularx}{\par}
\DeclareUnicodeCharacter{2192}{\ensuremath{\rightarrow}}
\definecolor{navy}{HTML}{15345B}
\definecolor{blue}{HTML}{1B5DE7}
\definecolor{teal}{HTML}{007F83}
\definecolor{ink}{HTML}{223248}
\definecolor{muted}{HTML}{5B6C82}
\definecolor{pale}{HTML}{EFF5FF}
\definecolor{mint}{HTML}{EDF8F5}
\definecolor{amber}{HTML}{8B560B}
\definecolor{sand}{HTML}{FFF7E9}
\hypersetup{pdftitle={Guía de formación desde cero: BI + IA aplicada a ventas},pdfauthor={Proyecto BI asistida por IA},pdfsubject={Formación, operación, demostración y escenarios para AdventureWorks y WideWorldImporters},pdfkeywords={BI, IA, formación, datamart, ventas, demostración, tesis, SQL Server}}
\color{ink}
\setlength{\parindent}{0pt}
\setlength{\parskip}{6pt}
\setlength{\emergencystretch}{2em}
\renewcommand{\arraystretch}{1.25}
\setlist[itemize]{leftmargin=5mm,itemsep=3pt,topsep=3pt,parsep=0pt}
\setlist[enumerate]{leftmargin=6mm,itemsep=4pt,topsep=4pt,parsep=0pt}
\newcolumntype{Y}{>{\raggedright\arraybackslash}X}
\pagestyle{fancy}
\fancyhf{}
\fancyhead[L]{\small\color{muted}BI + IA\enspace /\enspace Formación desde cero}
\fancyhead[R]{\small\color{muted}Ventas · SQL Server}
\fancyfoot[L]{\small\color{muted}Corte de referencia: 3 de octubre de 2026}
\fancyfoot[R]{\small\color{muted}\thepage\ /\ \pageref{LastPage}}
\renewcommand{\headrulewidth}{0.3pt}
\renewcommand{\footrulewidth}{0pt}
\newcommand{\eyebrow}[1]{{\small\bfseries\color{teal}\MakeUppercase{#1}}\par}
\newcommand{\pagetitle}[2]{\phantomsection\addcontentsline{toc}{section}{#2}\eyebrow{#1}\vspace{2mm}{\LARGE\bfseries\color{navy}#2\par}\vspace{3mm}}
\newcommand{\sub}[1]{\par\vspace{2mm}{\large\bfseries\color{navy}#1\par}}
\makeatletter
\renewcommand*\l@section[2]{\@dottedtocline{1}{0pt}{0pt}{#1}{#2}}
\makeatother
\newcommand{\code}[1]{{\small\texttt{\detokenize{#1}}}}
\newtcolorbox{goal}{enhanced,colback=pale,colframe=blue,boxrule=0.5pt,arc=2mm,left=4mm,right=4mm,top=3mm,bottom=3mm,title={Necesidad lista para copiar},fonttitle=\bfseries,colbacktitle=blue,coltitle=white}
\newtcolorbox{note}[1][Límite que debe explicar]{colback=sand,colframe=amber!25,boxrule=0.5pt,arc=2mm,left=3mm,right=3mm,top=2mm,bottom=2mm,title={#1},colbacktitle=sand,coltitle=amber,fonttitle=\bfseries}
\newtcolorbox{good}[1][Criterio de aceptación]{colback=mint,colframe=teal!25,boxrule=0.5pt,arc=2mm,left=3mm,right=3mm,top=2mm,bottom=2mm,title={#1},colbacktitle=mint,coltitle=teal,fonttitle=\bfseries}
\newcommand{\params}[5]{\sub{Parametrización recomendada}\begin{tabularx}{\linewidth}{@{}>{\bfseries}p{30mm}Y@{}}Fuente / dominio&#1 / Datamart de ventas\\Catálogo&#2\\Periodicidad&#3\\Grano a conservar&#4\\Dimensiones&#5\\\end{tabularx}}
\newcommand{\casepage}[4]{\clearpage\pagetitle{#1\enspace ·\enspace #2}{#3}{\small\color{muted}#4\par}\vspace{2mm}}
\newcommand{\copilot}[1]{\sub{Pregunta para el copiloto}\emph{«#1»}}
\newcommand{\evidence}[1]{\vfill{\footnotesize\color{muted}\textbf{Evidencia y alcance.} #1\par}}
\newcommand{\shot}[3]{\begin{center}\includegraphics[width=\linewidth,height=#3,keepaspectratio]{figuras/#1.png}\par\vspace{1mm}{\footnotesize\color{muted}Captura real de la auditoría del 3 de octubre de 2026. #2}\end{center}}
\newcommand{\chapterpage}[2]{\clearpage\pagetitle{#1}{#2}}

\begin{document}
\thispagestyle{empty}
\eyebrow{Manual de aprendizaje, operación y demostración}
\vspace{8mm}
{\fontsize{33}{37}\selectfont\bfseries\color{navy}BI + IA\par desde cero\par}
\vspace{4mm}
{\fontsize{23}{27}\selectfont\bfseries\color{navy}Aprender a transformar\par una necesidad de ventas\par en información comprobable\par}
\vspace{5mm}
{\Large\color{teal}AdventureWorks2022 + WideWorldImporters\par}
\vspace{5mm}
{\large Conceptos de negocio, uso paso a paso, capturas reales, doce necesidades de práctica y criterios para decidir cuándo confiar en el resultado.\par}
\vspace{5mm}
\begin{good}[La idea central que debe poder explicar]
La IA es un \textbf{asistente}: ayuda a formular, interpretar y proponer. La plataforma comprueba estructuras y ejecuta cálculos controlados. El analista decide el significado, revisa los límites y aprueba. Ni una redacción convincente ni un gráfico atractivo sustituyen la evidencia.
\end{good}
\sub{Para quién está escrita}
Para quien empieza sin conocimientos de BI y necesita usar la herramienta, estudiar el proyecto, preparar una exposición y reconocer errores de interpretación comercial. No se requiere escribir SQL ni programar para seguir el recorrido de usuario.
\begin{note}[Qué significa universal en esta guía]
Las reglas de aprobación, trazabilidad y cálculo son comunes. Las tablas y columnas que aquí se citan son referencias observadas de cada fuente, no una personalización que deba codificarse en el sistema. La evidencia disponible cubre estas dos bases SQL Server, no todos los motores o modelos posibles.
\end{note}
\vfill
{\small\color{muted}Versión 1.0 · 3 de octubre de 2026\par
Elaborada con el código, documentación y evidencia operativa de esta instalación. Capturas históricas reales identificadas como tales. No se repitieron cargas ni llamadas al LLM para redactar esta guía.}

\clearpage
\renewcommand{\contentsname}{Ruta de lectura}
{\small\tableofcontents}
\chapterpage{Módulo 1 · Orientación}{Cómo estudiar y qué podrá hacer}
\sub{Resultado de aprendizaje}
Al terminar podrá explicar una necesidad comercial, reconocer qué datos la respaldan, recorrer las pantallas sin confundir aprobación con ejecución, revisar un datamart y defender sus cifras frente a otra persona. No necesita memorizar nombres de tablas: necesita saber \textbf{qué preguntar y qué comprobar}.
\begin{tabularx}{\linewidth}{@{}p{31mm}YY@{}}
\toprule
\textbf{Recorrido}&\textbf{Qué estudiar}&\textbf{Qué producir}\\
\midrule
Primera lectura&Conceptos BI, proceso de ventas y papel de la IA.&Una explicación propia de pedido, factura, fila y KPI.\\
Práctica guiada&Uso paso a paso y una necesidad sencilla por fuente.&Ficha con objetivo, alcance, propuesta y evidencia.\\
Demostración&Guion breve, expedientes probados y preguntas al copiloto.&Una exposición reproducible, con límites explícitos.\\
Profundización&Doce escenarios, ejercicios y controles negativos.&Casos nuevos identificados como pendientes de ensayo.\\
\bottomrule
\end{tabularx}
\sub{Tres papeles que debe practicar}
\textbf{Analista BI:} traduce una decisión en datos y cálculos; elige el evento, revisa el grano y las relaciones, interpreta advertencias y documenta su aprobación.

\textbf{Gerente comercial:} pregunta qué ocurre, dónde, con qué magnitud y qué falta para decidir. No necesita aprobar fórmulas técnicas, pero sí entender qué mide cada cifra y qué no demuestra.

\textbf{QA o responsable de calidad:} repite pasos, compara esperado contra obtenido, registra fuente y filtros, prueba errores y evita afirmar que algo está probado sin evidencia.
\begin{note}[Perspectiva de trabajo no significa permiso informático]
En la auditoría se crearon los roles Analista BI y Gerente comercial, pero quedaron sin usuarios asignados. El recorrido integral se hizo con una cuenta administrativa autorizada. Cambiar entre Vista ejecutiva y Vista analítica no inicia sesión con otro rol ni prueba por sí mismo la seguridad de acceso.
\end{note}
\sub{Qué cubre y qué no cubre esta formación}
Cubre el dominio de ventas, dos fuentes SQL Server, catálogo por fuente, sugerencias de necesidades, validación, modelado supervisado, ETL, conciliación, analítica, exportación y copiloto. No enseña a ejecutar ventas en un ERP ni certifica contabilidad, tributación, cobranzas o costos históricos. No se implementa un pronóstico de ventas en estos ejercicios.
\sub{Cómo reconocer el nivel de evidencia}
\textbf{Probado integralmente:} el caso identificado llegó a ETL, panel y comprobaciones. \textbf{Previabilidad:} sólo se comprobó si podía continuarse. \textbf{Propuesto:} deberá ensayarlo. \textbf{Condicional:} depende de una definición o dato aún por confirmar. Una captura prueba lo que muestra, no todo lo que podría ocurrir después.
\begin{good}[Antes de pulsar cualquier aprobación]
Debe poder completar esta frase: «Estoy analizando \underline{\hspace{18mm}}, una fila representa \underline{\hspace{18mm}}, el importe significa \underline{\hspace{18mm}} y el resultado no permite afirmar \underline{\hspace{18mm}}».
\end{good}

\chapterpage{Módulo 1 · Conceptos de BI}{De registros a decisiones}
\textbf{BI} significa inteligencia de negocio: organizar datos y medidas para responder preguntas de gestión. No consiste sólo en dibujar gráficos. Primero se define qué ocurrió en el negocio; después se calcula y comunica con trazabilidad.
\begin{tabularx}{\linewidth}{@{}p{31mm}Y@{}}
\toprule
\textbf{Concepto}&\textbf{Explicación desde cero}\\
\midrule
Fuente OLTP&Base operativa que registra pedidos, facturas, productos o clientes. Está organizada para las transacciones del negocio.\\
Metadatos&Descripción de la estructura: tablas, columnas, tipos, claves y relaciones. No son los importes ni los nombres de todos los clientes.\\
Instantánea&Versión guardada de esos metadatos. Permite saber con qué estructura se construyó una propuesta. No es una copia de todas las ventas.\\
Datamart&Conjunto analítico para un área, aquí ventas. Reorganiza datos para consultar medidas por dimensiones.\\
Hecho y grano&El hecho representa el evento medido. El grano dice qué representa exactamente una fila, por ejemplo una línea de factura.\\
Dimensión&Conjunto descriptivo para clasificar hechos. Producto puede contener nombre, categoría y color. Esos campos son atributos de la dimensión, no medidas.\\
Medida&Valor sobre el que se calcula: cantidad, importe o costo. No todo número es una medida: un identificador no se suma.\\
KPI&Indicador definido para una pregunta: total facturado, número de pedidos o importe promedio por factura. Requiere fórmula y unidad.\\
ETL&Extracción, transformación y carga. Lee la fuente autorizada, aplica reglas controladas y guarda las tablas analíticas.\\
Linaje&Cadena que explica de dónde salió un resultado: fuente, instantánea, propuesta, ejecución, columna y cálculo.\\
\bottomrule
\end{tabularx}
\sub{Claves y relaciones, sin necesidad de programar}
Una \textbf{clave primaria (PK)} identifica una fila. Una \textbf{clave foránea (FK)} relaciona una tabla con otra. Una línea puede referirse a su factura; la factura puede referirse a su cliente. El sistema debe comprobar esa ruta. No basta con que dos columnas tengan nombres parecidos.

Una relación uno a muchos mal recorrida puede repetir líneas y aumentar ventas artificialmente. Por eso la aplicación revisa rutas, claves compuestas y conservación del grano. Un resultado que «parece razonable» no sustituye esa comprobación.
\textbf{Modelo estrella:} un hecho central se relaciona con dimensiones descriptivas. La propuesta usa la estructura que puede validar; no debe forzarse el mismo conjunto de tablas para todas las fuentes. Una métrica se convierte en instrumento de gestión al tener objetivo y contexto: esta guía no presupone que ya estén implementados metas, semáforos o alertas comerciales.
\begin{good}[Una diferencia decisiva]
Una factura con tres líneas representa \textbf{una factura, tres filas de detalle y una cantidad de unidades que puede ser distinta de tres}. Para contar facturas se usa el identificador distinto del documento, no el número de filas.
\end{good}

\chapterpage{Módulo 1 · El negocio de ventas}{Pedido, entrega, factura y cobro}
Una operación comercial puede atravesar varios eventos. No todas las empresas los realizan en el mismo orden ni cada documento corresponde uno a uno con otro. La necesidad debe nombrar el evento que se quiere medir.
\begin{tabularx}{\linewidth}{@{}p{27mm}YY@{}}
\toprule
\textbf{Evento}&\textbf{Qué indica}&\textbf{Qué no permite asumir}\\
\midrule
Pedido&Solicitud comercial registrada; productos, cantidades y condiciones.&Que se haya entregado, facturado o cobrado.\\
Entrega&Salida o entrega de bienes, si existe evidencia logística.&Que se haya cobrado o facturado todo.\\
Factura&Documento emitido con líneas e importes.&Que el dinero haya ingresado o que todo su total sea ingreso sin impuestos.\\
Cobro&Pago aplicado o entrada de dinero según registros verificables.&Que coincida con el mes de la factura o con una única factura.\\
Devolución / nota de crédito&Ajuste de la operación, si está representado y su relación es comprobable.&Que un importe negativo sea siempre una devolución o que el sistema la reste sin una regla.\\
\bottomrule
\end{tabularx}
\sub{Seis decisiones comerciales antes de pedir el datamart}
\begin{enumerate}
\item \textbf{Evento:} ¿pedidos registrados o facturas emitidas? «Ventas» por sí solo puede ser ambiguo.
\item \textbf{Población:} ¿todos los documentos o sólo un estado comprobado? No dé por hecho que se excluyeron anulados, borradores o devoluciones.
\item \textbf{Fecha:} ¿pedido, emisión, entrega o cobro? Para una factura mensual normalmente se revisa la fecha de emisión, no la de modificación de la fila.
\item \textbf{Importe:} ¿precio por cantidad, neto de descuento o total con impuesto? La palabra \emph{total} no responde esta pregunta.
\item \textbf{Cliente y territorio:} ¿comprador, cuenta de facturación, dirección de entrega o territorio comercial? Debe poder explicar la relación utilizada.
\item \textbf{Unidad y moneda:} ¿unidades físicas homogéneas? ¿una moneda comprobada? No sume indiscriminadamente monedas distintas ni convierta divisas sin una regla y una fuente.
\end{enumerate}
\begin{note}[Aplicación en nuestras fuentes]
AdventureWorks se demostró desde líneas de pedido de venta. WideWorldImporters tiene evidencia separada de pedidos y facturas. La diferencia entre sus totales \textbf{no demuestra crecimiento, error ni pendiente de cobro}: son poblaciones y definiciones distintas.
\end{note}
\sub{Periodicidad no es granularidad}
\textbf{Mensual} significa agrupar por año y mes. No agrupe sólo por «enero» si hay varios años. Conservar una fila por línea permite luego cambiar la agrupación sin perder trazabilidad. Tampoco significa que el programa programe una carga mensual automática.

\chapterpage{Módulo 1 · Ejemplo numérico de estudio}{Aprender a calcular antes de interpretar}
El ejemplo siguiente es \textbf{ficticio} y didáctico; no procede de las bases ni describe una regla tributaria. Se define un impuesto ilustrativo del 12\% sobre el neto y costos unitarios conocidos sólo para este ejercicio.
\begin{tabularx}{\linewidth}{@{}p{17mm}p{15mm}rrrrY@{}}
\toprule
\textbf{Doc.}&\textbf{Prod.}&\textbf{Cant.}&\textbf{Precio}&\textbf{Desc.}&\textbf{Neto}&\textbf{Costo total}\\
\midrule
P1&A&2&100&20&180&120\\
P1&B&1&50&0&50&30\\
P2&A&1&100&0&100&60\\
\midrule
Total&&4&&20&330&210\\
\bottomrule
\end{tabularx}
Hay \textbf{3 líneas, 2 documentos y 4 unidades}. El bruto antes de descuento es 350; el descuento monetario es 20; el neto antes de impuesto es 330. El impuesto del ejercicio es 39,60 y el total con impuesto es 369,60.
\begin{tabularx}{\linewidth}{@{}Yp{59mm}r@{}}
\toprule
\textbf{Indicador}&\textbf{Cálculo}&\textbf{Valor}\\
\midrule
Promedio neto por línea&330 / 3 líneas&110,00\\
Promedio neto por documento&330 / 2 documentos&165,00\\
Neto por unidad&330 / 4 unidades&82,50\\
Precio unitario promedio simple&(100 + 50 + 100) / 3&83,33\\
Margen bruto del ejercicio&330 - 210&120,00\\
Margen sobre venta neta&120 / 330 $\times$ 100&36,36\%\\
\bottomrule
\end{tabularx}
\sub{Qué le enseña este ejemplo}
\begin{itemize}
\item Los cuatro promedios contestan preguntas diferentes. No elija el de mayor valor: elija el que corresponde al nombre del indicador.
\item Una tasa de descuento no es dinero. Para una línea, precio por cantidad por tasa produce un importe de descuento, si ésa es la regla comprobada.
\item No reste el descuento dos veces de un importe que ya es neto. No mezcle importe con impuesto y costo sin impuesto para llamar al resultado margen comercial comparable.
\item El margen bruto no es utilidad neta: el ejemplo no incluye personal, alquiler, intereses ni otros gastos.
\item El margen porcentual del conjunto se calcula con totales compatibles. No se obtiene promediando porcentajes de líneas sin ponderación.
\end{itemize}
\begin{good}[Ejercicio resuelto]
Si filtra únicamente el producto B, quedan 50 de importe y un documento P1. El promedio filtrado es 50 por documento que contiene B; \textbf{no es el total completo de P1}, que en este ejemplo vale 230 antes de impuesto. Todo promedio debe explicar el alcance del filtro.
\end{good}

\chapterpage{Módulo 1 · Criterio comercial}{Cómo convertir una cifra en una conclusión}
\begin{tabularx}{\linewidth}{@{}p{31mm}YY@{}}
\toprule
\textbf{Pregunta de gestión}&\textbf{Qué revisar}&\textbf{Conclusión que debe evitar}\\
\midrule
¿Subieron las ventas?&Mismo evento, moneda, calendario y cobertura; períodos completos comparables.&«Crecimos» al comparar un mes parcial con otro completo.\\
¿Qué producto lidera?&Suma por producto; aclarar si el ranking es por importe o unidades.&Confundir la línea más cara con el producto de mayor contribución.\\
¿Qué cliente aporta más?&Identidad comercial, suma, documentos y alcance de cliente/cuenta.&Llamar fidelidad o abandono a un simple ranking de ventas.\\
¿Qué territorio funciona mejor?&Nivel geográfico y ruta: comercial, entrega, estado o país.&Atribuir causalidad al vendedor sin datos de responsabilidad comercial.\\
¿Qué tan rentable es?&Base neta, costo compatible y vigencia histórica comprobada.&Presentar costo estándar o último costo como costo real histórico.\\
¿Cuál es la participación?&Numerador y denominador: universo filtrado o sólo Top N.&Usar porcentaje del Top 5 como cuota de todas las ventas.\\
\bottomrule
\end{tabularx}
\sub{Chequeos operativos que también son de negocio}
\textbf{Cobertura:} un cliente sin nombre puede seguir contando en los importes. Revise cuántas identidades se resolvieron y cuáles siguen técnicas. No convierta un nulo automáticamente en cero ni lo elimine sin explicar el efecto.

\textbf{Comparabilidad:} cambiar una fuente, un expediente o una dimensión puede cambiar población, grano e importe. Lea siempre la cabecera antes de interpretar una tarjeta.

\textbf{Disponibilidad frente a calidad:} tener una columna de costo demuestra que existe; no demuestra que sea correcta, actualizada o histórica. Un metadato no certifica toda la realidad comercial.

\textbf{Causalidad:} el panel puede mostrar asociación, diferencia o concentración. Para explicar \emph{por qué} ocurrió algo se necesita evidencia adicional. Ni la IA ni una línea ascendente prueban una causa.
\begin{note}[Vocabulario seguro para la exposición]
Diga «en este expediente y con estos filtros se observa...», «el dato permite comparar...», «la causa no está demostrada...» o «el margen usa un costo de referencia...». Evite «la IA comprobó que el clima causó la caída» o «todo está correcto porque pasó el ETL».
\end{note}
\sub{Una necesidad bien formulada}
Indique \textbf{decisión + evento + medidas + dimensiones + fecha/período + grano + límites}. Ejemplo: «Para priorizar productos, comparar mensualmente importe y unidades de líneas de factura por producto, conservar factura y línea y no inferir cobros». Las tablas las propone el sistema; el significado lo debe revisar usted.

\chapterpage{Módulo 1 · Modelado dimensional}{Cómo se organiza un análisis de ventas}
Un \textbf{data warehouse} integra información analítica de un alcance amplio; un \textbf{datamart} se enfoca en un área, como ventas. Este proyecto materializa datamarts. No debe presentarse como una implantación completa de gobierno corporativo de datos para toda una empresa.
\sub{Un modelo estrella explicado con una venta}
Imagine una línea que vende dos unidades de un producto a un cliente en una fecha. El hecho conserva las medidas y las referencias a dimensiones; las dimensiones contienen descripciones para agrupar.
\begin{tabularx}{\linewidth}{@{}p{33mm}YY@{}}
\toprule
\textbf{Pieza}&\textbf{Qué contiene conceptualmente}&\textbf{Pregunta que permite responder}\\
\midrule
Hecho Venta&Referencia a fecha, producto, cliente y territorio; cantidad e importe; trazabilidad documental.&¿Cuánto ocurrió en cada línea y cómo se relaciona con el documento?\\
Dimensión Fecha&Día, mes y año de una fecha comercial definida.&¿Cuánto se registró por período?\\
Dimensión Producto&Identificador, nombre y atributos verificables.&¿Qué productos aportaron más?\\
Dimensión Cliente&Identificador comercial y descripción comprobada.&¿A quién corresponde el importe?\\
Dimensión Territorio&Identidad y nivel territorial declarado.&¿Dónde se atribuye comercial o geográficamente?\\
\bottomrule
\end{tabularx}
La tabla anterior es un modelo didáctico, no una obligación de generar exactamente cinco tablas en cualquier necesidad. Categoría y color pueden ser atributos de Producto; no todo atributo requiere una tabla nueva. Las relaciones reales de la fuente condicionan lo que se puede construir.
\sub{Por qué se conserva el detalle}
Si guarda sólo el total de cada mes, ya no puede verificar qué factura contribuyó ni separar correctamente productos después. Con un grano de línea, puede agregar por mes, producto o territorio conservando la trazabilidad. No debe unir cabeceras repetidas y sumar de nuevo su total por cada línea: duplicaría dinero.
\sub{Medidas aditivas y no aditivas}
\textbf{Aditivas:} cantidades e importes compatibles pueden sumarse a través de grupos sin duplicación. \textbf{No aditivas:} porcentajes, precios medios y ratios no se suman como importes. Un conteo distinto tampoco suele poder sumarse entre productos: el mismo pedido puede contener varios productos.
\begin{good}[Ejemplo de doble conteo]
P1 contiene A y B. El conteo por A es 1 y por B es 1; sumarlos produce 2, pero sigue existiendo \textbf{un solo pedido}. El total documental debe volver a calcular el identificador distinto sobre el universo total.
\end{good}

\chapterpage{Módulo 1 · Calidad y análisis}{De dato disponible a evidencia útil}
\sub{Seis preguntas sobre calidad de datos}
\begin{tabularx}{\linewidth}{@{}p{33mm}Y@{}}
\toprule
\textbf{Dimensión de calidad}&\textbf{Cómo pensarla en este proyecto}\\
\midrule
Completitud&¿Hay importe, fecha e identidad donde se necesitan? Un nulo debe tener tratamiento explícito; no convertirlo silenciosamente en cero.\\
Unicidad&¿La clave identifica realmente la fila? ¿Una unión está repitiendo el hecho?\\
Validez&¿El tipo y la operación son compatibles? Una fecha o un identificador no son dinero.\\
Consistencia&¿Panel, expediente, exportación y copiloto describen el mismo alcance cuando se comparan?\\
Exactitud&¿La cifra representa lo que ocurrió? La conciliación aporta evidencia, pero no demuestra que toda captura original del negocio sea correcta.\\
Oportunidad&¿El corte de datos sirve para la decisión? Una carga hecha hoy puede contener ventas históricas y no actualización en tiempo real.\\
\bottomrule
\end{tabularx}
La plataforma aporta controles de estas características, pero no se debe afirmar que calcula automáticamente un índice completo de calidad para todas ellas. El analista debe revisar también las limitaciones del proceso de negocio.
\sub{Describir, diagnosticar y predecir son tareas distintas}
\textbf{Descriptivo:} qué importe, unidades o documentos se observan. \textbf{Comparativo:} diferencias entre grupos o períodos equivalentes. \textbf{Diagnóstico:} explicar causas requiere evidencia adicional, no sólo una narrativa. \textbf{Predictivo:} estimar futuro requiere un modelo y evaluación específica; no es el alcance implementado de estas demostraciones.
\sub{Operaciones analíticas que aprenderá}
\textbf{Agregar:} sumar por mes. \textbf{Filtrar:} restringir a un año o territorio. \textbf{Desglosar:} comparar categorías de un gráfico. \textbf{Cambiar nivel:} pasar conceptualmente de día a mes o de ciudad a país sólo si el contrato y la interfaz lo soportan. No se promete un explorador OLAP libre con cualquier jerarquía.
\begin{note}[Métrica, KPI y meta]
Una métrica es una medida calculada. Un KPI se elige porque importa para una decisión. Para decir «se cumplió el objetivo» además necesita una meta, plazo y criterio de comparación. El hecho de que la interfaz llame KPI al total de ventas no significa que configure automáticamente cuotas, alertas o semáforos de desempeño.
\end{note}
\sub{Trazabilidad y reproducibilidad}
Conservar entrada, fuente, estructura, contrato, fórmula y ejecución permite explicar un resultado. Reproducir la validación de un contrato no significa que el LLM vaya a redactar la misma respuesta cada vez. La evidencia estable pertenece al artefacto validado y a su ejecución, no a una promesa de infalibilidad del modelo.

\chapterpage{Módulo 2 · IA como asistente}{Qué hace un LLM y por qué se le valida}
Un \textbf{LLM} es un modelo que genera e interpreta lenguaje. Puede traducir una necesidad a una propuesta estructurada y redactar una explicación. Su capacidad lingüística no le convierte en dueño de las reglas de su empresa ni garantiza que cada columna o fórmula sea correcta.
\sub{Palabras que verá al usarlo}
\begin{tabularx}{\linewidth}{@{}p{31mm}Y@{}}
\toprule
\textbf{Término}&\textbf{Qué significa para el usuario}\\
\midrule
Proveedor / modelo&Servicio y modelo configurados. Cambiarlo puede variar latencia y calidad de respuesta, pero no debe cambiar los controles de aprobación.\\
Contexto&Información que se entrega para una tarea: necesidad, estructura permitida o agregados autorizados, según la etapa. No equivale a acceso irrestricto a la base.\\
Prompt / instrucción&Petición y reglas que orientan la respuesta. No sustituye la comprobación posterior.\\
Token / presupuesto&Unidad de procesamiento y límite de salida. Una respuesta cortada puede fallar; aumentar razonamiento no garantiza corrección.\\
Salida estructurada&Respuesta con un formato contractual, normalmente JSON, que el software puede comprobar. No es un archivo que el analista deba editar.\\
Alucinación&Referencia, fórmula o afirmación plausible pero no respaldada. El sistema debe descartarla o dejarla por revisar.\\
Regla determinística&Comprobación reproducible basada en código y datos, no en una nueva opinión del modelo. Por ejemplo, si una columna existe o si un conteo coincide.\\
\bottomrule
\end{tabularx}
\sub{Cuatro niveles distintos de confianza}
\begin{enumerate}
\item \textbf{Formato:} la respuesta se puede leer y cumple el contrato.
\item \textbf{Estructura:} las tablas, tipos, claves y relaciones existen y son compatibles.
\item \textbf{Semántica:} representan lo solicitado: costo no es precio; factura no es pedido; tasa no es dinero. Hay reglas y también límites que requieren revisión humana.
\item \textbf{Datos:} la ejecución y sus agregados se concilian con la fuente. Esto ocurre después, no al redactar la necesidad.
\end{enumerate}
\begin{good}[La defensa correcta del proyecto]
El valor de la IA es asistir una interpretación compleja y reducir trabajo manual. El valor del sistema es \textbf{acotar, comprobar, conservar evidencia y exigir decisiones humanas}. Rechazar una respuesta inválida es una protección; no debe ocultarse ni convertirse en aprobación automática.
\end{good}

\chapterpage{Módulo 2 · Responsabilidades}{Quién hace qué en cada etapa}
{\small
\begin{tabularx}{\linewidth}{@{}p{24mm}YYY@{}}
\toprule
\textbf{Etapa}&\textbf{IA}&\textbf{Plataforma}&\textbf{Persona}\\
\midrule
Conexión y metadatos&No inventa ni captura el esquema.&Prueba acceso y obtiene estructura del motor; conserva instantáneas.&Autoriza fuente y revisa contexto.\\
Catálogo&No es la autoridad que publica preguntas por sí sola.&Mantiene preguntas por fuente y descubre evidencia técnica.&Define objetivos y preguntas de negocio.\\
Sugerir / redactar necesidad&Propone textos con el contexto permitido.&Valida referencias y capacidades; separa lo adoptable de lo bloqueado.&Elige o mantiene el texto y vuelve a validar.\\
Viabilidad&Interpreta requisitos, posibilidades y límites.&Combina revisión con reglas, tipos, rutas y huella del contexto.&Revisa cada límite y acepta sólo un alcance entendido.\\
Conceptos y propuesta&Propone conceptos, dimensiones, medidas, KPI y plan.&Normaliza y valida contrato y cobertura; rechaza referencias inválidas.&Comprueba evento, grano y significado.\\
Personalización&Puede asistir una decisión contextual cuando se solicita.&Ofrece selecciones y recetas controladas; valida versiones.&Incluye/excluye, elige opciones válidas y justifica cambios.\\
Verificar / aprobar&No concede la aprobación.&Recalcula reglas y evidencia; registra decisión.&Aprueba o rechaza explícitamente.\\
Materializar / ETL&No ejecuta SQL libre del modelo.&Construye y ejecuta el plan controlado; carga y concilia.&Confirma una ejecución y revisa sus resultados.\\
Etiquetas descriptivas&Puede sugerir normalización o traducción elegible.&Conserva identificadores y limita categorías enviables.&Revisa y publica etiquetas cuando corresponda.\\
Panel y hallazgos&Los hallazgos automáticos del panel no requieren LLM.&Calcula KPI, gráficos y resúmenes determinísticos.&Elige filtros, compara y evalúa significado.\\
Copiloto&Interpreta consulta acotada y redacta respuesta.&Obtiene agregados autorizados, impone límites y conserva contexto.&Contrasta cifras y cuestiona conclusiones.\\
Exportación&No redacta ni recalcula libremente el informe.&Construye PDF/Excel desde la selección autorizada.&Verifica alcance y uso del documento.\\
\bottomrule
\end{tabularx}}
\begin{note}[No todo lo automático es inteligencia artificial]
Sumar ventas, filtrar un año, detectar una columna inexistente y comparar conteos son tareas del software. Una explicación en español puede provenir de reglas o del LLM; compruebe qué componente está viendo antes de describirlo como «lo hizo la IA».
\end{note}

\chapterpage{Módulo 2 · Seguridad y límites}{Qué información se usa y qué no se garantiza}
\sub{La información depende de la etapa}
\begin{itemize}
\item \textbf{Necesidad, conceptos y propuesta:} texto del objetivo y metadatos permitidos, como nombres de tablas y columnas, tipos y relaciones. No se envían credenciales ni filas de ventas para estos pasos.
\item \textbf{Etiquetas durante el flujo de materialización:} pueden enviarse categorías no sensibles de baja cardinalidad para sugerir descripciones. No debe afirmarse que jamás sale ningún valor de datos en todo el sistema.
\item \textbf{Copiloto analítico:} recibe contexto y agregados autorizados para responder. Un total comercial sigue siendo información de negocio; su organización debe autorizar ese uso.
\item \textbf{Proveedor local:} una configuración local cambia el destino de procesamiento, no elimina la necesidad de validar respuestas, permisos y alcance.
\end{itemize}
\sub{Qué significa independencia del modelo y la fuente}
El contrato, las reglas y la trazabilidad son comunes. La conexión y la instantánea determinan qué existe en cada base; el adaptador del proveedor determina cómo se solicita una respuesta. No se debería resolver un fallo escribiendo reglas exclusivas para una base de ejemplo o aceptando fórmulas que otro modelo rechazó.

Eso \textbf{no equivale a compatibilidad ya demostrada con cualquier motor, idioma, esquema o modelo}. Esta guía está respaldada por dos fuentes SQL Server. Las pruebas sintéticas en español e inglés amplían la cobertura, pero no sustituyen ensayos de una nueva instalación.
\sub{Límites conocidos que debe poder explicar}
\begin{itemize}
\item Una sugerencia puede quedar bloqueada o no ser útil. En la última repetición de reformulación AW se mantuvo el texto original; no se obtuvo una reformulación utilizable.
\item Groq encontró un límite temporal de cuota en la auditoría. Claude permitió completar el caso facturado después de correcciones. No se deduce de eso una clasificación estadística de proveedores.
\item Una columna existente no certifica moneda, impuesto, estado comercial o costo histórico. «Comprobado estructuralmente» no significa «certificado financieramente».
\item Roles creados sin usuarios asignados no prueban segregación de funciones. Esa prueba autenticada sigue pendiente en la evidencia citada.
\end{itemize}
\begin{good}[Regla de seguridad para aprender]
Trabaje con la fuente de demostración, no pegue secretos en la necesidad ni en el chat, no cambie permisos para sortear un bloqueo y no acepte límites sólo para llegar a la siguiente pantalla. Si el sistema no puede resolver algo, documente qué falta.
\end{good}

\chapterpage{Módulo 3 · Primer acceso}{Entrar y orientarse sin conocimientos previos}
\sub{Antes de abrir la herramienta}
El responsable técnico debe haber iniciado la aplicación, restaurado las dos bases y entregado una cuenta autorizada. \textbf{Registrar una conexión no instala SQL Server ni restaura una base inexistente}. Si la página no abre, no intente resolverlo modificando datos: solicite que comprueben los servicios. La instalación y los comandos de infraestructura pertenecen al manual técnico del proyecto.
\begin{enumerate}
\item Abra el navegador. En esta instalación de demostración la dirección es \url{http://localhost:5173/}; en otro equipo use la dirección indicada por su administrador. \emph{Localhost} se refiere al propio equipo donde se abre el navegador.
\item Complete \textbf{Correo} y \textbf{Contraseña}. Pulse \textbf{Iniciar sesión}. No proyecte ni copie la contraseña en la guía, en una necesidad o en una captura.
\item Compruebe el usuario en el encabezado. Abra \textbf{Menú} si la barra lateral no se ve; el diseño cambia con el ancho de la ventana.
\item Entre en \textbf{Inicio}. Si su perfil dispone de la preparación, revise fuente, metadatos y asistente de IA. «Todo está listo» sólo indica prerrequisitos. El usuario final puede ir directamente a \textbf{Analítica de ventas} si su permiso no incluye diseño.
\item Localice \textbf{Fuente en contexto} y \textbf{Cambiar fuente}. Elija la base del ejercicio y compruebe su nombre antes de continuar.
\end{enumerate}
\sub{El menú como mapa de trabajo}
\begin{tabularx}{\linewidth}{@{}p{49mm}Y@{}}
\toprule
\textbf{Grupo}&\textbf{Para qué se utiliza}\\
\midrule
Preparación del entorno&Conexiones de datos, Explorador de esquema, Configuración LLM y Parámetros. Prepara las condiciones de trabajo.\\
Diseño y transformación BI&Catálogo analítico, Asistente de datamart y Datamart de ventas. Define, revisa y materializa el análisis.\\
Análisis y decisiones&Analítica de ventas. Consulta resultados, gráficos, hallazgos, conversación y exportaciones.\\
Administración y seguridad&Usuarios, roles y permisos disponibles según la cuenta. No es una etapa que el alumno deba modificar para avanzar.\\
\bottomrule
\end{tabularx}
\begin{note}[Si no aparece una opción]
Puede depender del permiso de su cuenta. Pida al administrador el acceso que corresponde a su función. No use una cuenta ajena ni confunda una opción oculta con un fallo de datos. El perfil gerencial debería consultar y exportar; no necesita administrar credenciales para ver el panel.
\end{note}
\sub{Primer ejercicio sin riesgo}
Cambie de AdventureWorks a WideWorldImporters y vuelva a Inicio. Diga en voz alta qué fuente está seleccionada. Esta acción cambia su contexto de trabajo; no deshabilita ni elimina la otra base.

\chapterpage{Módulo 3 · Preparación autorizada}{Conectar la fuente y dejar lista la IA}
\sub{A. Conexiones de datos}
Si AdventureWorks y WideWorldImporters ya están registradas, \textbf{no las duplique}. Revise su nombre y estado. Para una fuente nueva, el responsable autorizado completa:
\begin{tabularx}{\linewidth}{@{}p{43mm}Y@{}}
\toprule
\textbf{Campo}&\textbf{Qué debe introducir}\\
\midrule
Nombre visible&Nombre que permita reconocer negocio y entorno, sin incluir contraseñas.\\
Motor&Actualmente SQL Server; no se ofrece aquí soporte operativo de cualquier motor.\\
Servidor / Puerto&Dirección y puerto entregados por el administrador, accesibles desde el servicio de la aplicación.\\
Base de datos&Nombre exacto de la base restaurada. No basta con cambiar el nombre visible.\\
Usuario de sólo lectura&Cuenta con permisos mínimos para consultar la fuente y sus metadatos.\\
Contraseña&Secreto de esa cuenta. No se escribe en objetivos ni en comentarios.\\
Transporte&Opciones de cifrado y certificado revisadas por el administrador. No desactive seguridad para resolver un error sin diagnóstico.\\
\bottomrule
\end{tabularx}
Pulse \textbf{Registrar conexión}, luego \textbf{Probar conexión}. Sólo cuando se valide el acceso de sólo lectura, use \textbf{Habilitar}. Después pulse \textbf{Leer metadatos}. Las fuentes pueden permanecer habilitadas simultáneamente.
\sub{B. Configuración LLM}
Con una configuración activa y probada, el alumno no necesita crear otra. Si debe configurarse una: entre en \textbf{Configuración LLM}, complete nombre, proveedor, URL, modelo y razonamiento; use \textbf{Crear registro}. Para cloud: \textbf{Registrar credencial} y \textbf{Guardar credencial}; después \textbf{Probar conexión} y \textbf{Activar}.

Cambiar proveedor, modelo o nivel exige volver a comprobar la configuración. Use exclusivamente las opciones admitidas por el formulario y por su cuenta. En el corte documentado, Claude ofrece Automático, Mínimo y Bajo; no se debe imponer el mismo nivel a todos los proveedores.
\begin{note}[Dos cambios distintos]
\textbf{Cambiar fuente} selecciona la base de trabajo. \textbf{Activar una configuración LLM} cambia el proveedor activo para las operaciones posteriores y puede afectar a otras sesiones. No lo cambie durante una exposición compartida sin coordinarlo.
\end{note}
\begin{good}[Punto de control]
Puede continuar cuando reconoce la fuente, el acceso está probado, existe una instantánea y el proveedor está listo. Una prueba LLM correcta sólo acredita conexión y respuesta básica; no garantiza todas las propuestas posteriores.
\end{good}

\chapterpage{Módulo 3 · Actividad 1}{Leer el esquema antes de pedir un análisis}
Ruta: \textbf{Preparación del entorno → Explorador de esquema}.
\begin{enumerate}
\item Compruebe \textbf{Fuente en contexto}. Si no hay captura, use \textbf{Crear primera instantánea}; si ya existe, consúltela antes de decidir si necesita \textbf{Actualizar metadatos}.
\item Lea versión, fecha, huella y conteos. \textbf{Versión de metadatos} permite consultar capturas guardadas; con una sola versión puede estar deshabilitada con una explicación.
\item En \textbf{Buscar tabla o columna}, busque una entidad del caso, pulse \textbf{Buscar} y seleccione la tabla. Revise columnas, tipos, PK y relaciones entrantes/salientes.
\item Identifique detalle, cabecera, fecha comercial, producto y cliente. En WWI distinga líneas de pedido de líneas de factura.
\end{enumerate}
\shot{01_explorador_wwi}{Fuente WWI e instantánea \#3; el explorador consulta estructura guardada.}{112mm}
\begin{good}[Lo que debe aprender a decir]
«Existe una línea de factura relacionada con su cabecera; allí está la fecha de emisión». Es una afirmación estructural. «Todas las facturas están cobradas» no puede deducirse de esa relación.
\end{good}
\textbf{Importante:} consultar una instantánea histórica en el Explorador no la convierte en selección global para una nueva propuesta. El asistente utiliza la instantánea vigente de la fuente. No debe editar los metadatos para justificar una respuesta del modelo.
Actualizar metadatos \textbf{no recarga las ventas}. El panel consulta los datos de su ejecución; su «Última actualización» no significa que los documentos comerciales sean de esa fecha.

\chapterpage{Módulo 3 · Actividad 2}{Entender y conservar el catálogo de ventas}
Ruta: \textbf{Diseño y transformación BI → Catálogo analítico}.
\begin{enumerate}
\item Compruebe fuente y \textbf{Dominio habilitado: Datamart de ventas}.
\item Revise \textbf{Cobertura técnica de la fuente} y sus referencias detectadas. «Verificable» no significa que se haya cargado un datamart; «Por validar» no es una evidencia definitiva.
\item Lea las \textbf{Preguntas de negocio} y las \textbf{Periodicidades}. Para el primer ejercicio no cambie el catálogo.
\end{enumerate}
\shot{02_catalogo_aw}{Catálogo propio de AdventureWorks y separación entre pregunta comercial y descubrimiento técnico.}{108mm}
\sub{¿Quién construye el catálogo?}
El sistema parte de un perfil de ventas y descubre evidencia desde los metadatos. El analista configura las preguntas; \textbf{no se publica automáticamente un catálogo particular completo sólo porque lo diga la IA}. En esta instalación AW conserva seis preguntas y WWI cuatro.

Para una modificación autorizada: \textbf{Agregar pregunta}, completar Etiqueta, Explicación para el analista e Instrucción funcional para la IA; después \textbf{Guardar catálogo}. Añadir «¿por qué cayó la venta por clima?» no crea una fuente meteorológica.
\begin{note}[No confunda restaurar con refrescar]
\textbf{Restaurar catálogo validado} puede reemplazar la personalización. No lo use durante el ejercicio para actualizar la pantalla, especialmente si ya hay preguntas propias de la empresa.
\end{note}

\chapterpage{Módulo 3 · Actividad 3}{Escribir una necesidad o pedir sugerencias}
Ruta: \textbf{Asistente de datamart → Datamart de ventas → Crear propuesta}.
\begin{enumerate}
\item Escriba en \textbf{Objetivo del análisis} un texto de una ficha: entre 20 y 2.000 caracteres para validar. No escriba SQL ni secretos.
\item Marque al menos una de las \textbf{Preguntas de negocio disponibles} pertinentes y seleccione \textbf{Periodicidad: Mensual}.
\item Lea el \textbf{Destino} proveedor/modelo. Marque la autorización para enviar esa necesidad y los metadatos permitidos a ese destino.
\item Pulse \textbf{Validar viabilidad}, o solicite ayuda con las opciones siguientes si necesita formular primero el objetivo.
\end{enumerate}
\shot{33_necesidad_metadatos_consentimiento}{Campo inicial, consentimiento y preguntas WWI. La captura muestra las acciones de IA bloqueadas hasta autorizar.}{88mm}
\sub{Tres caminos válidos para el principiante}
\textbf{Ya sé qué necesito:} escriba el objetivo, configure preguntas/período, autorice y valide.

\textbf{No sé por dónde empezar:} autorice y pulse \textbf{Sugerir necesidades con esta fuente}; lea las opciones y adopte una sólo si representa su decisión comercial.

\textbf{Tengo una idea mal redactada:} use \textbf{Ayúdame a formular la necesidad}. Si la alternativa está habilitada y le sirve, pulse \textbf{Usar esta redacción}; si no, \textbf{Mantener mi redacción}.
\begin{note}[El consentimiento se revisa al cambiar el objetivo]
Adoptar una sugerencia o cambiar texto, preguntas, período o contexto invalida la evaluación y puede desmarcar el consentimiento. Termine de editar, vuelva a autorizar el destino y valide de nuevo. No es un fallo que le pida revisar el envío actualizado.
\end{note}

\chapterpage{Módulo 3 · Actividad 4}{Interpretar la viabilidad y los derivables}
Después de \textbf{Validar viabilidad}, lea cada requisito. Abra \textbf{Ver evidencia técnica}; no mire únicamente el total de elementos verdes.
\begin{tabularx}{\linewidth}{@{}p{38mm}Y@{}}
\toprule
\textbf{Resultado}&\textbf{Qué hacer como analista}\\
\midrule
Directo / referencia comprobada&Confirmar que la referencia representa el concepto solicitado, no sólo que existe.\\
Derivable / automático&Revisar columnas y relación o receta candidata. La propuesta posterior debe implementarla y validarla; usted no tiene que escribir SQL.\\
Ambiguo / por decidir&Leer el límite. Si no cambia la decisión principal, aceptarlo conscientemente; si la impide, editar el objetivo.\\
No disponible&No inventar una fuente. Reducir explícitamente el alcance o conseguir evidencia adicional antes de continuar.\\
\bottomrule
\end{tabularx}
\shot{36_necesidad_wwi_revision}{Caso NEED-18: estructura suficiente y límites pendientes. La continuación no es todavía una carga ETL.}{106mm}
\begin{note}[No marque todas las advertencias por rutina]
Si quiere rentabilidad histórica y sólo hay último costo, aceptar esa limitación implica \textbf{otro alcance}, no resolver mágicamente el dato histórico. Si la limitación invalida la decisión, deténgase y reformule.
\end{note}
Cuando las condiciones se cumplen se habilita \textbf{Generar propuesta y resolver N derivables}. Al pulsarlo empieza el diseño asistido. Puede consumir IA y producir un rechazo seguro; ni la viabilidad ni ese botón garantizan todos los KPI.

\chapterpage{Módulo 3 · Actividad 5}{Separar interpretación de fórmula ejecutable}
La interfaz distingue la interpretación libre del modelo de una operación que el sistema puede materializar. Esta separación evita tomar una frase convincente como si ya fuera un cálculo ejecutado.
\shot{38_necesidad_aw_interpretacion}{La etiqueta «Interpretación propuesta por IA (no ejecutable)» es deliberada. Las referencias comprobadas no certifican por sí solas la composición financiera del importe.}{119mm}
\sub{Cómo leer una tarjeta del ejemplo}
\begin{enumerate}
\item \textbf{Nombre:} «Importe registrado de línea de venta» es una descripción prudente. No presupone que sea bruto, neto o cobrado.
\item \textbf{Interpretación:} una suma propuesta expresa la intención. Aún debe traducirse al contrato controlado de la propuesta.
\item \textbf{Referencia:} al abrir la evidencia, compruebe tabla y columna de la fuente correcta.
\item \textbf{Advertencia:} si la composición del importe no está probada, conserve el nombre neutral o documente evidencia adicional antes de renombrarlo.
\end{enumerate}
\begin{good}[Comprobación oral para la demostración]
Pregunte: «¿La existencia de \code{UnitPrice} demuestra un costo?». Respuesta: no; puede ser precio de venta. Tampoco una fecha o un identificador numérico pueden convertirse en medida monetaria sólo porque el modelo lo proponga.
\end{good}

\chapterpage{Módulo 3 · Actividades 6 y 7}{Revisar conceptos y diseño del datamart}
\sub{Paso 2 del asistente: Conceptos}
Revise qué conceptos aparecen como \textbf{Incluir}. Los de evidencia débil pueden venir excluidos. Abra \textbf{Ver evidencia y cómo resolver}; el \textbf{Consultar al copiloto} contextual es opcional y sirve para entender una decisión, no para aprobarla automáticamente.

No use \textbf{Confirmar inclusión excepcional} sólo para avanzar. Una conversación o recomendación no cambia sola la selección: revise y aplique únicamente una opción comprobable. Continúe con \textbf{Generar propuesta BI}. Si cambia conceptos puede generarse otra versión y consumir IA; si no cambió la selección, el botón puede limitarse a continuar.
\sub{Paso 3: Propuesta}
\begin{tabularx}{\linewidth}{@{}p{41mm}Y@{}}
\toprule
\textbf{Sección que revisa}&\textbf{Pregunta que debe responder}\\
\midrule
Granularidad&¿Una fila es una línea de pedido o factura? ¿Se conserva su identificador?\\
Hecho y medidas&¿Se eligió el evento correcto? ¿Cada importe tiene fuente, unidad y agregación?\\
Dimensiones&¿Fecha, producto, cliente y territorio llegan mediante relaciones verificadas sin multiplicar las líneas?\\
KPIs propuestos&¿Suma, conteo, promedio o cociente responden a la pregunta? ¿Se identifica el denominador?\\
De la necesidad al datamart&¿Se implementó cada requisito o se muestra su límite? ¿Quedaron sólo los operandos, pero falta el KPI solicitado?\\
Plan ETL declarativo&¿Describe transformaciones permitidas? No es SQL libre que el analista deba copiar y ejecutar.\\
Validaciones y advertencias&¿Hay errores bloqueantes? ¿Las exclusiones o límites están comprendidos?\\
\bottomrule
\end{tabularx}
\sub{Comprobar nombres y cobertura}
Este botón revisa identidades descriptivas en la fuente dentro de la plataforma. Un identificador de cliente puede ser válido aunque el nombre aún no se haya resuelto. La comprobación no es otra opinión del LLM ni significa que se haya ejecutado el ETL completo.
\begin{good}[Lista mínima antes de seguir]
Evento correcto; una fila bien definida; fecha comercial; importe con significado explícito; documento contado de forma distinta; dimensiones sin duplicar; requisitos cubiertos o limitados de forma visible.
\end{good}
\textbf{Error típico:} una propuesta de facturación usa líneas de pedidos. No se arregla cambiando el título a «facturas». Debe corregirse el origen del hecho y sus relaciones, o rechazarse.

\chapterpage{Módulo 3 · Actividad 8}{Personalizar sin romper lo aprobado}
Use \textbf{Personalizar propuesta} cuando exista una razón de negocio. Si está conforme, \textbf{Continuar a revisión} o \textbf{Continuar sin cambios}. Personalizar no debe ser un paso obligatorio para hacer una demostración vistosa.
\sub{Qué se puede revisar en el formulario}
\begin{itemize}
\item Incluir o excluir dimensiones, medidas e indicadores que el contrato permite.
\item Elegir una agregación compatible: Suma, Conteo, Conteo distinto, Promedio, Mínimo o Máximo.
\item Revisar la \textbf{Medida compatible} de un KPI y, cuando corresponda, \textbf{Definir cálculo controlado} con columnas y operaciones permitidas.
\item Registrar \textbf{Justificación del ajuste} y pulsar \textbf{Guardar como nueva versión}. Se conserva la procedencia de la versión anterior.
\end{itemize}
\sub{Guía para elegir una agregación}
\begin{tabularx}{\linewidth}{@{}p{30mm}YY@{}}
\toprule
\textbf{Opción}&\textbf{Uso razonable}&\textbf{Precaución}\\
\midrule
Suma&Unidades o importes aditivos compatibles.&No sumar identificadores, tasas o monedas incompatibles.\\
Conteo&Número de filas o valores, según contrato.&No llamarlo documentos si el hecho es una línea.\\
Conteo distinto&Pedidos o facturas únicas.&Verificar la clave documental y el filtro.\\
Promedio&Media simple de los valores de fila.&No es automáticamente precio ponderado o ticket documental.\\
Mínimo / Máximo&Menor o mayor valor de una población.&Un máximo de línea no identifica al producto líder por suma.\\
\bottomrule
\end{tabularx}
\textbf{Cocientes agregados:} no son otra opción del selector Agregación. El KPI importe total / documentos distintos debe quedar correctamente propuesto y compilado. La operación \textbf{Dividir} de un cálculo por fila no sustituye ese cociente de agregados; no la use para fabricarlo manualmente.
\sub{Ejemplo de una justificación útil}
«Se excluye el indicador de producto líder calculado como máximo de línea, porque la necesidad requiere ordenar la suma por producto. Se conserva el importe total para el ranking analítico y se comprueba que la exclusión no deje requisitos esenciales sin cubrir».
\begin{note}[No es edición libre de SQL]
Cambiar sólo el texto del grano no cambia su implementación física. La corrección guiada de relaciones requiere evidencia y es una tarea avanzada. Si no comprende sus consecuencias, conserve la versión y pida revisión; no fuerce una unión ni habilite un KPI por su nombre.
\end{note}
\textbf{Criterio:} más KPI no significa mejor datamart. Un indicador excluido con una razón verificable puede ser una decisión de calidad superior a una cifra incorrecta.

\chapterpage{Módulo 3 · Actividad 9}{Verificar la evidencia y aprobar conscientemente}
\textbf{Verificar evidencia} revisa metadatos, referencias, contrato y reproducción determinística. No llama al LLM, no lee las filas comerciales para conciliar importes y no ejecuta ETL.
\begin{enumerate}
\item Compruebe objetivo, fuente, periodicidad y número de versión.
\item Confirme que no haya errores bloqueantes. Lea el detalle, no sólo el color.
\item Lea las advertencias y acepte únicamente las que comprende y que no contradicen su decisión comercial.
\item Escriba un \textbf{Comentario de revisión} con evento, grano, fórmula crítica y límites aceptados.
\item Pulse \textbf{Aprobar propuesta}; si no corresponde, \textbf{Rechazar propuesta} con motivo o \textbf{Personalizar antes de decidir}.
\end{enumerate}
\shot{08_validacion_aw}{Revisión supervisada de AW. La propia pantalla separa aprobación, comprobación estructural y ejecución futura.}{108mm}
\begin{good}[Ejemplo de comentario]
«Revisadas fuente y fecha comercial. Se conserva una fila por línea y el número de pedidos usa documentos distintos. El promedio se calcula con importe total dividido por pedidos distintos. Acepto los límites descritos; no interpreto el resultado como cobros».
\end{good}
\textbf{Resultado esperado:} una versión aprobada queda disponible para la materialización compatible. \textbf{Aprobar no crea tablas ni carga datos}. Consultar o verificar una versión histórica tampoco debe retirar su aprobación por sí solo.

\chapterpage{Módulo 3 · Actividad 10}{Preparar y ejecutar el ETL}
Ruta: \textbf{Diseño y transformación BI → Datamart de ventas}.
\begin{enumerate}
\item Si se abre un expediente existente, use \textbf{Volver a propuestas} para iniciar una ejecución nueva. Compruebe que eligió su versión aprobada, no sólo la primera tarjeta.
\item \textbf{Seleccionar} → \textbf{Revisar indicadores}. Lea fórmula, unidad, cobertura y KPI seleccionados.
\item Pulse \textbf{Revisar transformaciones}. Examine plan, advertencias y decisiones; escriba el \textbf{Registro de decisión} y marque su confirmación.
\item Pulse \textbf{Preparar materialización}. Esto prepara un expediente; aún no es la carga.
\item Compruebe fuente y versión otra vez. Pulse \textbf{Materializar y cargar datamart} una sola vez y espere el resultado.
\end{enumerate}
\shot{12_conciliacion_aw}{Las cinco etapas del espacio de materialización. Esta captura muestra la llegada al paso de validación, no los conteos de conciliación.}{101mm}
\sub{Qué ocurre detrás de la pantalla}
La plataforma lee las columnas autorizadas de SQL Server, aplica operaciones permitidas y guarda hecho y dimensiones en PostgreSQL. El destino está versionado por ejecución; una carga no debe mezclar AdventureWorks con WideWorldImporters. No se ejecuta SQL libre devuelto por el LLM.
\begin{note}[Acción con efecto real]
Materializar crea datos analíticos. No lo pulse repetidamente para «ver si avanza». Si falla, abra el expediente y revise el mensaje; no borre datamarts ni retire aprobaciones para reiniciar la demostración.
\end{note}

\chapterpage{Módulo 3 · Actividad 11}{Conciliar, revisar etiquetas y reconocer la moneda}
Después de cargar, revise \textbf{Conciliación OLTP--datamart}: filas de origen, filas cargadas, diferencia, calidad de las dimensiones e indicadores conciliados. Un conteo de filas igual es necesario en estos casos, pero no prueba por sí solo todas las definiciones de negocio.
\shot{31_wwi_facturas_etl15}{WWI, expediente \#15: 228.265 líneas en origen y destino, diferencia cero y moneda todavía sin código comprobado.}{124mm}
\sub{Qué significa cada observación}
La dimensión Fecha puede tener menos filas que el hecho: muchas líneas comparten una fecha. «Duplicados controlados» en una dimensión no implica automáticamente duplicación de ventas. Lo importante es conservar una única identidad dimensional compatible y no multiplicar el hecho.

\chapterpage{Módulo 3 · Actividad 11, continuación}{Resolver moneda y etiquetas antes de publicar}
Puede haber datos cargados y conciliados, pero quedar una decisión de publicación o interpretación. Son estados diferentes; no repita la carga para resolver un texto o una divisa.
\sub{Si la divisa no está comprobada}
\begin{enumerate}
\item Lea el aviso y la unidad de los indicadores. Un total monetario calculable no acredita una moneda única.
\item Si aparece \textbf{Comprobar divisa sin repetir el ETL}, úselo con el permiso correspondiente y revise la evidencia obtenida.
\item Si sigue sin existir un código único comprobado, conserve \textbf{moneda de origen}. No seleccione USD por similitud con otro expediente ni convierta valores sin tasa y regla autorizadas.
\item En una comparación o informe, explique la limitación. Puede describir el total registrado; no combinarlo con importes de otra moneda como si fueran equivalentes.
\end{enumerate}
\sub{Si existen etiquetas candidatas pendientes}
\begin{enumerate}
\item Revise valor original y etiqueta sugerida. Una traducción cambia la presentación, no debe cambiar identidad ni nivel de la categoría.
\item Corrija o desmarque lo que no corresponda. No traduzca identificadores, nombres personales o texto libre como si fueran categorías comerciales.
\item Registre comentario, confirme la revisión y use \textbf{Publicar etiquetas revisadas}. Este botón aparece cuando existen mapeos que requieren esa decisión; no es un botón genérico para crear otro datamart.
\item Si falló únicamente el proveedor al interpretar etiquetas, use \textbf{Reintentar interpretación sin repetir el ETL}, cuando esté disponible. Conserve la conciliación anterior.
\end{enumerate}
\begin{note}[Ejemplo que debe entender el alumno]
Europe puede presentarse como Europa si el mapeo es correcto y conserva la identidad. Eso no transforma un grupo territorial en un país o estado. La traducción de una etiqueta no añade detalle geográfico.
\end{note}
\begin{good}[Puede pasar a analítica cuando...]
La ejecución está validada/conciliada y publicada, con datos físicos consultables. Si aparece Interpretar o Revisar, resuelva lo señalado. Con diferencias bloqueantes, deténgase y conserve la evidencia. No renombre un error como «dato no importante».
\end{good}

\chapterpage{Módulo 3 · Actividad 12}{Leer el panel y probar los gráficos}
Ruta: \textbf{Análisis y decisiones → Analítica de ventas}.
\begin{enumerate}
\item Lea fuente, ejecución y propuesta en la cabecera. No interprete tarjetas antes de hacerlo.
\item Con varias ejecuciones consultables, elija \textbf{Datamart analizado}. Con una sola puede mostrarse una identidad fija; no es necesario un selector vacío.
\item Empiece por \textbf{Vista ejecutiva}; pruebe \textbf{Año} y \textbf{Territorio}. Use años presentes en los datos, no el año actual por costumbre.
\item Pulse un punto, barra o categoría. Lea su valor y \textbf{Participación visible}; \textbf{Cerrar detalle} sólo cierra la ficha. \textbf{Filtrar tablero por...}, cuando existe, sí cambia el filtro.
\item En \textbf{Vista analítica}, seleccione \textbf{Indicador de los gráficos} y abra \textbf{Consultar los datos de este gráfico}. Después pulse \textbf{Restablecer filtros}.
\end{enumerate}
\shot{14_dashboard_aw}{Identidad AW: propuesta \#98 y ejecución \#14. Los modos ejecutivo y analítico son vistas, no un cambio de usuario.}{107mm}
\begin{note}[Tres lecturas que no son equivalentes]
Un valor nulo o «Sin valor» no es cero. Un porcentaje del conjunto visible no siempre usa el total de la empresa. Una variación entre extremos no prueba una tendencia sostenida ni explica su causa; puede incluir meses parciales.
\end{note}
Los \textbf{Hallazgos explicables} del panel se calculan mediante reglas. Lea \textbf{Cómo se generaron}; no los confunda con una conversación nueva con el LLM.

\chapterpage{Módulo 3 · Actividad 13}{Conversar con el copiloto y verificarlo}
En \textbf{Preguntar sobre estos resultados}, escriba una pregunta concreta y pulse \textbf{Preguntar al copiloto}, o use una sugerencia guiada. Esta acción sí puede consumir IA.
\sub{Cinco puntos que debe revisar en cada respuesta}
\textbf{Evento:} pedido o factura. \textbf{Indicador:} importe, cantidad o conteo. \textbf{Filtros:} año y territorio realmente consultados. \textbf{Denominador:} documentos, unidades o total del universo. \textbf{Límite:} qué no permite concluir.
\shot{32_wwi_facturas_copiloto2015}{Consulta real WWI en 2015: 62.090.220,81, 22.250 facturas y promedio 2.790,57. La respuesta distingue facturación de pedidos y cobros.}{121mm}
\sub{Preguntas útiles para aprender}
«¿Qué documento estás contando y cómo calculas su promedio?». «¿Qué cinco productos aportan más y cuál es el denominador de su participación?». «¿Qué limitaciones debo considerar antes de tomar una decisión?».
\begin{note}[Una pregunta puede cambiar el alcance de la consulta]
El copiloto puede consultar un filtro permitido que no esté seleccionado en el tablero. Eso no debe confundirse con un cambio automático de los gráficos. Lea la consulta interpretada y el alcance declarado antes de comparar cifras. Si afirma una causa no respaldada, señálela y no la use como evidencia gerencial.
\end{note}

\chapterpage{Módulo 3 · Lectura guiada}{Leer una respuesta con Top N y participación}
En esta conversación de AdventureWorks se preguntó por cinco productos en Southwest. Sirve para estudiar cómo debe declararse el universo de una participación; no para aceptar cada redondeo sin comprobarlo.
\shot{15_copiloto_aw}{Conversación histórica AW. La respuesta identifica territorio, productos e importe; la consulta interpretada muestra el Top 5.}{125mm}
\sub{Cómo auditar la explicación}
\begin{enumerate}
\item Identifique \textbf{indicador} y \textbf{orden}: importe por producto de mayor a menor, no una línea máxima.
\item Compruebe que los cinco productos pertenecen al territorio consultado y que el período está declarado.
\item Localice el \textbf{denominador}. Para cuota territorial debe usarse el total territorial, no solamente la suma de los cinco productos.
\item Compare valores detallados con la evidencia agregada. La auditoría dejó señalada la revisión del redondeo narrativo de centavos: una explicación correcta del alcance no certifica cada centavo escrito por el modelo.
\end{enumerate}
\begin{note}[Regla para el usuario final]
Si un porcentaje no coincide, antes de concluir que hay un error compruebe universo y redondeo. Si todavía difiere, registre el caso y use la cifra verificable; no corrija el número mentalmente mientras comunica una conclusión.
\end{note}

\chapterpage{Módulo 3 · Actividad 14}{Cambiar de fuente sin mezclar eventos}
Repita el ejercicio con la otra base. El propósito es demostrar que el mismo procedimiento funciona con estructuras diferentes, no hacer que ambas produzcan los mismos nombres o totales.
\begin{enumerate}
\item Anote la fuente y ejecución actual. Quite filtros que no desee conservar.
\item Use \textbf{Cambiar fuente}. Revise de nuevo nombre de base, instantánea, catálogo y expedientes.
\item En WWI distinga el expediente de \textbf{pedidos} del de \textbf{facturas}. Las tarjetas pueden ser válidas en ambos y contestar preguntas distintas.
\item Vuelva a AW y confirme su propia propuesta/ejecución. No deberían quedar productos, territorios o respuestas visibles de la fuente anterior como si pertenecieran a ésta.
\end{enumerate}
\shot{24_dashboard_wwi_alcance}{Expediente WWI \#13: la cabecera advierte que son pedidos registrados. No es el expediente de facturas \#15.}{112mm}
\begin{good}[Qué debe decir durante la demostración]
«Aquí cambió el contexto, no se sumaron bases. La fuente y el expediente gobiernan qué datos se consultan. En este historial hay pedidos; en el siguiente hay facturas. Voy a verificar el evento antes de comparar».
\end{good}
\textbf{Si detecta mezcla:} detenga la interpretación, anote fuente, versión, filtros y pasos; capture la pantalla y registre un incidente. No reconstruya o borre datamarts hasta distinguir un problema visual de uno de datos.

\chapterpage{Módulo 3 · Cierre de la operación}{Exportar, consultar históricos y empezar otra vez}
\sub{Exportar un resultado entendible}
\begin{enumerate}
\item En Analítica, deje seleccionados fuente, expediente, vista, indicador y filtros que quiere comunicar.
\item Pulse \textbf{Exportar PDF} o \textbf{Exportar Excel}, si su cuenta tiene permiso.
\item Abra el archivo descargado. Verifique procedencia, unidades, filtros y correspondencia con los gráficos visibles.
\item Acompañe el informe con una conclusión y sus límites. No distribuya información de negocio fuera de su audiencia autorizada.
\end{enumerate}
El PDF no es una captura del navegador: se construye desde la selección autorizada. El Excel contiene agregados que respaldan los gráficos visibles, no todas las filas operativas. No se debe prometer que la exportación incluya toda la conversación del copiloto.
\sub{Abrir una versión sin consumir de nuevo IA}
En el Asistente: \textbf{Versiones generadas → Mostrar → Aprobadas}, o el estado que necesite; pulse \textbf{Abrir resultado}. Las cinco etapas se consultan como historial. No se llama otra vez al modelo sólo por abrir una versión.

Para otra necesidad pulse \textbf{Crear nueva propuesta}. No cierre sesión para desbloquear el asistente ni use \textbf{Retirar aprobación} o \textbf{Descartar versión} como si fueran botones de regreso.
\sub{Consultar un expediente sin repetir el ETL}
En Datamart de ventas, busque \textbf{Expedientes recientes} y use \textbf{Abrir expediente \#...}. Recorrer sus etapas muestra evidencia; no repite la carga. Un expediente puede seguir siendo auditable aunque no tenga tablas físicas disponibles para analítica o no sea elegible para una nueva ejecución.
\begin{note}[Propuesta, revisión de necesidad y ejecución tienen números diferentes]
La revisión de viabilidad \#5 no es la propuesta \#5. Una propuesta \#114 puede originar la ejecución \#15. Registre el tipo de identificador junto al número. La versión recomendada para una nueva materialización no tiene por qué ser la que alimenta el panel que está consultando.
\end{note}
\sub{Si vence la sesión}
Vuelva a \textbf{Iniciar sesión} con su cuenta autorizada. Compruebe contexto y borrador restaurados antes de continuar. El consentimiento puede necesitar renovarse si cambió el objetivo o el destino. No repita una carga sólo porque el navegador le pidió autenticarse.
\begin{good}[Cierre de la sesión de práctica]
Guarde la ficha del caso y sus capturas, registre si fue consulta histórica o ejecución nueva y use \textbf{Cerrar sesión}. No documente contraseñas ni tokens en el acta.
\end{good}


\clearpage
\pagetitle{Preparación de los ejercicios}{Configuración común de las doce fichas}
\begin{enumerate}
\item \textbf{Seleccione la fuente en contexto.} Compruebe su nombre en Inicio, Explorador, Asistente y Analítica. No cambie credenciales ni retire versiones aprobadas para comenzar otra demostración.
\item \textbf{Revise el esquema.} Confirme la instantánea, tabla de detalle, fecha comercial, identificador del documento, claves y relaciones. Los conteos históricos de referencia son AW: 71 tablas y 90 relaciones; WWI: 48 tablas y 98 relaciones. No son requisitos fijos de otra instalación.
\item \textbf{Seleccione Datamart de ventas y Mensual.} En las fichas, mensual es la periodicidad del análisis; no programa automáticamente el ETL ni cambia una fila de detalle por una fila mensual. Puede ensayar otras periodicidades sólo si están habilitadas en su catálogo.
\item \textbf{Revise el proveedor.} Use una configuración que supere su prueba de conexión y un nivel de razonamiento admitido. El envío para formular o validar la necesidad contiene texto y estructura permitida, no filas ni credenciales.
\item \textbf{Escriba objetivo y preguntas; después autorice y valide.} Complete las entradas antes de marcar el consentimiento. Si adopta una sugerencia, vuelva a autorizar el contexto cambiado. Una sugerencia de IA es una opción, no una aprobación. Si la reformulación se bloquea, conserve su texto y vuelva a evaluarlo.
\end{enumerate}
\sub{Preguntas que debe marcar}
Las abreviaturas siguientes son sólo de esta guía; no son códigos para introducir en la aplicación. Las etiquetas corresponden al catálogo observado en esta instalación.
\begin{tabularx}{\linewidth}{@{}p{18mm}Y@{}}
\toprule
\textbf{Fuente}&\textbf{Etiquetas exactas y abreviaturas de lectura}\\
\midrule
AW&\textbf{T}: Evolución de ventas en el tiempo.\newline
\textbf{P}: Desempeño por producto y características.\newline
\textbf{C}: Desempeño por cliente.\newline
\textbf{G}: Desempeño por territorio.\newline
\textbf{R}: Costos y rentabilidad de las ventas.\newline
\textbf{U}: Ventas y costos por unidad.\\
\midrule
WWI&\textbf{T}: Evolución de ventas en el tiempo.\newline
\textbf{P}: Productos con mayor desempeño.\newline
\textbf{C}: Desempeño por cliente.\newline
\textbf{G}: Desempeño por territorio.\\
\bottomrule
\end{tabularx}
\begin{note}[No todos los parámetros son campos iniciales]
El formulario inicial recoge objetivo, preguntas y periodicidad. El grano, las dimensiones, las medidas y los KPI de las fichas son \textbf{criterios para revisar la propuesta y personalizar las opciones verificadas}, no un formulario libre para introducir SQL o forzar tablas. Una pregunta seleccionada orienta a la IA; no garantiza un KPI ni una fórmula concreta.
\end{note}
\sub{Cómo resolver lo que aparece como derivable}
Revise la ruta o fórmula candidata, sus columnas y su significado. Continúe para que la propuesta construya el cálculo controlado y lo valide. Si falta evidencia o se altera el grano, no lo fuerce: limite la necesidad o excluya ese componente con una decisión explícita. Aceptar una advertencia no convierte una interpretación incierta en un hecho comprobado.

\chapterpage{Módulo 4 · Banco de ejercicios}{Elegir una necesidad para cada demostración}
Todas las fichas conservan detalle documental y usan periodicidad Mensual. Los KPI son \textbf{criterios para revisar el contrato}, no campos que deba escribir libremente. Las letras del catálogo se explicaron en la página anterior.
\begin{tabularx}{\linewidth}{@{}p{18mm}Yp{41mm}r@{}}
\toprule
\textbf{Caso}&\textbf{Qué se aprende}&\textbf{Tipo de evidencia}&\textbf{Pág.}\\
\midrule
AW-01&Tendencia e importe por pedido&Capacidad E2E previa&\pageref{aw01}\\
AW-02&Objetivo sencillo por producto&Previabilidad del texto&\pageref{aw02}\\
AW-03&Cliente e identidad descriptiva&Nuevo texto propuesto&\pageref{aw03}\\
AW-04&Territorio y denominador&Capacidad analítica previa&\pageref{aw04}\\
AW-05&Importe ponderado por unidad&Nuevo texto propuesto&\pageref{aw05}\\
AW-06&Descuento, costo y margen&Avanzado condicional&\pageref{aw06}\\
\midrule
WWI-01&Facturación integral&Texto/alcance E2E previo&\pageref{wwi01}\\
WWI-02&Objetivo mínimo de factura&Previabilidad del texto&\pageref{wwi02}\\
WWI-03&Producto facturado y unidad&Capacidad E2E previa&\pageref{wwi03}\\
WWI-04&Cliente y ticket por factura&Nuevo texto propuesto&\pageref{wwi04}\\
WWI-05&Pedido no equivale a factura&Capacidad E2E histórica&\pageref{wwi05}\\
WWI-06&Geografía y trazabilidad&Nuevo texto propuesto&\pageref{wwi06}\\
\bottomrule
\end{tabularx}
\sub{Dos formas de usar el banco}
\textbf{Exposición estable:} abra los resultados conservados AW propuesta \#98 / ejecución \#14 y WWI propuesta \#114 / ejecución \#15, si siguen disponibles. Explique qué se ejecutó y muestre sus cifras. Es una consulta histórica, no una ejecución nueva en directo.

\textbf{Práctica de generación:} comience por AW-02 o WWI-02 para aprender viabilidad; después amplíe a un caso de negocio. Registre el número de la nueva versión y no espere que repita exactamente la redacción o las decisiones de otra respuesta del modelo.
\sub{Cómo registrar cada intento}
Anote fuente, instantánea, proveedor/modelo, texto, preguntas y período; luego qué límite apareció, cuál se aceptó y por qué. Si sólo habilitó la generación, escriba \textbf{previabilidad}. Si materializó, anote conciliación, ejecución y filtros usados para comprobarla.
\begin{note}[No restaure datos ni aprobaciones para igualar el ejemplo]
Los números de versión, ejecución y los valores de referencia pertenecen al corte auditado. Si su fuente cambió, un resultado diferente no implica por sí solo un fallo. Compare el mismo evento, población, fórmula y corte temporal, y explique cualquier diferencia.
\end{note}
\sub{Recomendación para la tesis}
No mezcle ensayos previos con nuevos. Identifique cada ejecución, conserve las capturas y describa incidentes y recuperación. No calcule una tasa de éxito de proveedores con intentos cuyo código, objetivo o configuración fueron cambiando.

\casepage{AW-01}{Básico}{Pedidos por mes y promedio documental}{Respaldado por capacidad probada integralmente; texto propuesto para repetir.}\label{aw01}
\begin{goal}
Comparar mensualmente el importe registrado de las líneas de venta, las unidades vendidas y el número de pedidos distintos. Calcular el importe promedio por pedido como el importe total dividido entre los pedidos distintos, conservando la trazabilidad de cada línea y sin asumir que los pedidos fueron cobrados.
\end{goal}
\params{AdventureWorks2022}{T: Evolución de ventas en el tiempo.}{Mensual}{Una línea de pedido de venta; conservar pedido y línea.}{Fecha comercial; producto, cliente o territorio sólo si se incluyen y verifican.}
\sub{Qué revisar en la propuesta}
\begin{itemize}
\item Importe y unidades con \textbf{Suma}; pedido con \textbf{Conteo distinto}.
\item Promedio por pedido = importe total / pedidos distintos bajo el mismo filtro.
\item Referencias observadas: \code{Sales.SalesOrderDetail} y cabecera \code{Sales.SalesOrderHeader}, con \code{OrderDate}. Son puntos de revisión, no instrucciones para imponer tablas.
\end{itemize}
\begin{good}
La ejecución histórica \#14, propuesta \#98, concilió 121.317 líneas y 31.465 pedidos. El importe fue 109.846.381,4250 y el promedio 3.491,07. Úselos como referencia sólo con el mismo corte, fórmula y población.
\end{good}
\copilot{¿Cuántos pedidos distintos hay y por qué el promedio por pedido no es el promedio por línea? Indica el numerador y el denominador.}
\begin{note}
\code{AVG(LineTotal)} es promedio de línea; contar todas las apariciones de un pedido cuenta líneas. No presente pedidos registrados como dinero cobrado.
\end{note}
\evidence{QA-03 a QA-07 del acta operativa. La funcionalidad se probó hasta panel y copiloto; el texto de esta ficha no se volvió a generar durante la preparación del manual.}

\casepage{AW-02}{Primer ensayo de IA}{Importe y unidades por producto}{Texto exacto que superó la revisión previa; no se ejecutó otro ETL en ese ensayo.}\label{aw02}
\begin{goal}
Comparar por mes el importe registrado y la cantidad de unidades de las líneas de venta por producto, conservando su identificador de pedido.
\end{goal}
\params{AdventureWorks2022}{P: Desempeño por producto y características.}{Mensual}{Una línea de venta, sin perder el pedido original.}{Fecha comercial y producto; atributos adicionales sólo si están verificados.}
\sub{Qué revisar en la propuesta}
\begin{itemize}
\item Suma del importe de línea y suma de cantidades; identificador de producto como clasificación, no como importe.
\item Relación del detalle con la cabecera para la fecha y con el producto para su identidad.
\item Si aparecen categoría, subcategoría u otros atributos, verificar su ruta y evitar que multipliquen el hecho. No todos los atributos se convierten automáticamente en filtros del panel.
\end{itemize}
\begin{good}[Qué ocurrió en la prueba previa]
La revisión de necesidad \#6 tuvo 11 directos, 1 derivable, 5 advertencias y 0 no disponibles. Aceptar los cinco límites habilitó continuar, pero \textbf{no se pulsó generar ni se cargó otro datamart}. Los conteos no representan 17 KPI únicos.
\end{good}
\copilot{¿Qué producto aporta más importe y cuál aporta más unidades? Si son distintos, explica qué mide cada clasificación sin inventar una causa.}
\begin{note}
El importe registrado no debe etiquetarse automáticamente como cobrado, sin impuesto o limitado a pedidos completados. Si una reformulación de IA se bloquea, mantenga este texto y vuelva a validar su alcance; no altere los metadatos para aceptar la sugerencia.
\end{note}
\evidence{NEED-17 y NEED-19. La reformulación final no produjo una alternativa utilizable; la previabilidad del texto original sí fue positiva. La pregunta al copiloto de esta ficha es una práctica propuesta para después de materializar.}

\casepage{AW-03}{Intermedio}{Clientes reconocibles y población comercial}{Nuevo texto propuesto; requiere una ejecución propia para declararlo aprobado.}\label{aw03}
\begin{goal}
Comparar mensualmente el importe registrado, las unidades vendidas y los pedidos distintos por cliente. Mostrar un nombre descriptivo y el tipo de cliente cuando puedan obtenerse mediante relaciones verificadas, conservando su identificador original y sin incluir personas ajenas a la población comercial.
\end{goal}
\params{AdventureWorks2022}{C: Desempeño por cliente; T: Evolución de ventas en el tiempo.}{Mensual}{Una línea de pedido; conservar cliente y pedido de origen.}{Fecha y cliente; nombre y tipo condicionados a su evidencia.}
\sub{Qué revisar en la propuesta}
\begin{itemize}
\item Importe/unidades con Suma y pedidos con Conteo distinto.
\item Cliente comercial relacionado desde la cabecera, no todas las personas de la base.
\item Enriquecimiento de \code{Sales.Customer} hacia persona o tienda sólo por rutas verificadas. Un cliente organización no tiene por qué tener persona individual.
\end{itemize}
\begin{good}
Use \textbf{Comprobar nombres y cobertura}. El conjunto de clientes debe mantener su población y sus identificadores. Después del ETL, los importes no deben aumentar por duplicaciones de la relación descriptiva.
\end{good}
\copilot{¿Cuáles son los principales clientes por importe dentro del período consultado? Aclara qué identidad representa cliente y qué porcentaje del total aporta cada uno.}
\begin{note}
Un nombre ausente requiere revisión, no eliminación automática de la venta. Un ranking no demuestra fidelidad, abandono o potencial de compra. No intercambie cliente comercial con persona, contacto o dirección.
\end{note}
\evidence{La resolución descriptiva y las dimensiones de cliente tienen evidencia en el proyecto, pero este objetivo concreto es un caso de práctica nuevo. La presencia del atributo en el contrato no garantiza un gráfico dedicado en todas las vistas.}

\casepage{AW-04}{Intermedio}{Territorios y contribución comercial}{Capacidad de filtros y Top N probada; objetivo compuesto nuevo.}\label{aw04}
\begin{goal}
Comparar mensualmente el importe registrado, las unidades vendidas y los pedidos distintos por territorio comercial. Identificar los territorios y productos con mayor contribución dentro del período seleccionado, indicando el conjunto usado como denominador de cada participación y respetando el nivel territorial existente en la fuente.
\end{goal}
\params{AdventureWorks2022}{G: Desempeño por territorio; P: Desempeño por producto y características; T: Evolución de ventas en el tiempo.}{Mensual}{Una línea de pedido, relacionada con su territorio comercial.}{Fecha, territorio comercial y producto.}
\sub{Qué revisar en la propuesta y el panel}
\begin{itemize}
\item Importe/unidades sumados y pedidos distintos. Territorio ligado al evento mediante relaciones comprobadas.
\item Producto líder calculado mediante suma agrupada por producto, no máximo de una sola línea.
\item Nivel territorial explícito: grupo, territorio, estado o país no son sinónimos. Una traducción de etiqueta no cambia su nivel.
\end{itemize}
\begin{good}[Prueba que puede mostrar]
En AW \#14 se ensayó 2014: 11.761 pedidos y promedio de 1.705,46. También se consultó Top 5 en Southwest con un denominador territorial de 24.184.609,60. No mezcle los filtros de ambos controles.
\end{good}
\copilot{¿Cuáles son los cinco productos con mayor importe en Southwest? Indica la participación sobre todo el territorio y diferencia ese denominador de la suma del Top 5.}
\begin{note}
Una participación visible en el gráfico puede usar sólo las categorías mostradas. Compare explícitamente los denominadores antes de afirmar que el panel y el copiloto se contradicen.
\end{note}
\evidence{QA-07 documenta filtro y conversación. La nueva necesidad de esta ficha requiere revisar de nuevo el contrato; no se promete que el modelo seleccione exactamente el mismo nivel de territorio.}

\casepage{AW-05}{Intermedio}{Venta por unidad y promedios ponderados}{Nuevo texto propuesto para aprender agregación y denominadores.}\label{aw05}
\begin{goal}
Comparar mensualmente el importe registrado por unidad vendida entre productos. Sumar primero el importe y las unidades de las líneas de venta y dividir el importe total entre las unidades totales dentro de cada grupo. Conservar el pedido de origen, mostrar la cantidad usada como denominador y no confundir este indicador con el promedio simple de precios unitarios. No incluir costos en este ejercicio.
\end{goal}
\params{AdventureWorks2022}{P: Desempeño por producto y características; T: Evolución de ventas en el tiempo.}{Mensual}{Una línea de venta con cantidad e importe compatibles.}{Producto y fecha.}
\sub{Qué revisar en la propuesta}
\begin{itemize}
\item Dos medidas base: importe y cantidad, ambas sumadas.
\item Importe por unidad = suma de importes / suma de cantidades. Revisar tratamiento de cantidad nula o cero.
\item Si la IA propone promedio simple del precio, exigir otra etiqueta o una fórmula compatible con la necesidad; no aceptarlo como sustituto silencioso.
\end{itemize}
\begin{good}
Compare el total de importe y el total de unidades del mismo alcance. Su cociente debe coincidir, dentro del redondeo mostrado, con el KPI por unidad. En cada agregado por producto, ambos operandos deben corresponder al mismo producto; no se presupone un filtro de producto en el tablero.
\end{good}
\copilot{Explica la diferencia entre importe por unidad, precio unitario promedio simple y promedio por pedido. ¿Cuál de ellos está calculando este expediente?}
\begin{note}
No se marca U sólo para añadir costos: esta necesidad excluye costos deliberadamente. Si el catálogo orienta la respuesta a un alcance mayor, revise y mantenga el objetivo. Un valor por unidad no es automáticamente un precio de lista ni una utilidad.
\end{note}
\evidence{Las recetas de cociente de agregados tienen pruebas en el proyecto; este texto y su nuevo resultado deben ensayarse. Sirve para reconocer el error observado anteriormente de promedios etiquetados de forma incorrecta.}

\casepage{AW-06}{Avanzado condicional}{Descuentos, costos y margen de referencia}{Aprobar sólo con definición comercial y evidencia compatibles.}\label{aw06}
\begin{goal}
Comparar mensualmente por producto el importe registrado, las unidades y el descuento monetario calculable con referencias verificadas. Incorporar costo total, venta por unidad, costo por unidad y margen sólo si existe una fuente de costo semánticamente comprobada. Si es un costo estándar de referencia y no el costo histórico de la venta, indicar esa limitación y no presentar el margen como rentabilidad histórica real.
\end{goal}
\params{AdventureWorks2022}{R: Costos y rentabilidad de las ventas; U: Ventas y costos por unidad; P: Desempeño por producto y características.}{Mensual}{Una línea de venta; operandos trazables y relacionados sin duplicación.}{Producto y fecha, ampliaciones sólo verificadas.}
\sub{Qué revisar en la propuesta}
\begin{itemize}
\item Descuento monetario: precio por tasa por cantidad cuando ésa sea la receta comprobada; no sumar tasas como dinero.
\item Costo total: cantidad por costo unitario válido. \code{Production.Product.StandardCost} es candidato de costo de referencia, no prueba automática de costo histórico.
\item Venta/costo por unidad: cociente de totales. Margen: base de venta compatible menos costo; porcentaje con denominador definido.
\end{itemize}
\begin{note}
No convierta \code{UnitPrice} en costo. No reste nuevamente descuentos incluidos en el importe. Si falta evidencia, excluya el KPI o limite explícitamente el alcance; no declare cubierta la rentabilidad histórica solicitada con una simple estimación.
\end{note}
\copilot{¿Este margen usa costo histórico o un costo de referencia? Explica fórmula, base del importe y qué decisión no debería tomarse con esta limitación.}
\evidence{Existen recetas financieras y pruebas históricas, entre ellas \#91/\#12. La ficha es avanzada y condicional: una conciliación numérica no certifica la interpretación histórica del costo ni convierte el margen en utilidad neta.}

\casepage{WWI-01}{Demostración integral}{Facturación emitida para la gerencia}{Objetivo persistido de la propuesta \#114; ciclo integral probado en \#15.}\label{wwi01}
\begin{goal}
Para priorizar productos y territorios, analizar mensualmente ventas efectivamente facturadas a nivel de línea de factura. El importe comercial es el total de la línea de factura con impuesto tal como consta en la fuente, no el precio unitario por cantidad antes de impuesto. Comparar ese importe y las unidades por producto, cliente y territorio; contar facturas distintas y calcular importe promedio por factura como total facturado dividido para facturas distintas. Usar la fecha de emisión de la factura, conservar el identificador de línea y factura y el pedido de origen cuando la relación esté verificada. Sólo usar atributos y relaciones comprobados.
\end{goal}
\params{WideWorldImporters}{T, P, C y G: las cuatro preguntas del catálogo WWI.}{Mensual}{Una línea de factura; conservar factura y línea.}{Fecha de emisión, producto, cliente y territorio verificado.}
\sub{Qué revisar y mostrar}
\begin{itemize}
\item Hecho de factura, no de pedido; referencia observada \code{Sales.InvoiceLines}, cabecera \code{Sales.Invoices} y \code{InvoiceDate}.
\item Importe/unidades sumados; facturas distintas; ticket = total / facturas distintas.
\item En este corte se comprobó que \code{ExtendedPrice} incluye impuesto: 172.261.341,20 de precio por cantidad más 25.782.098,25 de impuesto.
\end{itemize}
\begin{good}
El expediente \#15 concilió 228.265 líneas, 70.510 facturas y 198.043.439,45. Promedio por factura: 2.808,73; sin código de moneda único demostrado. Use «moneda de origen».
\end{good}
\copilot{Para 2015, ¿cuál es el importe total facturado, cuántas facturas distintas hubo y cuál fue el promedio por factura? Explica el denominador y si son pedidos, facturas o cobros.}
\evidence{QA-14 a QA-16. La composición fiscal se confirmó en esta instalación por conciliación; no se infiere sólo del nombre ExtendedPrice ni se traslada a otras fuentes. Repetir hoy la generación es un ensayo nuevo.}

\casepage{WWI-02}{Primer ensayo de IA}{Una necesidad mínima de facturación}{Texto exacto de previabilidad NEED-18; no una nueva carga probada.}\label{wwi02}
\begin{goal}
Analizar mensualmente el importe registrado de las líneas de factura, conservando la trazabilidad con su factura y sin inferir cobros ni composición de impuestos.
\end{goal}
\params{WideWorldImporters}{T: Evolución de ventas en el tiempo.}{Mensual}{Una línea de factura, conservando factura y línea.}{Fecha comercial de la factura; otras dimensiones sólo si se incluyen y verifican.}
\sub{Qué revisar en la propuesta}
\begin{itemize}
\item Importe de línea con Suma; fecha de emisión y vínculo con la factura.
\item No aceptar una tabla de pedidos como sustituto de facturas; no usar fecha de modificación como fecha comercial por defecto.
\item Conservar el nombre neutral «importe registrado» si no se aporta la evidencia de composición que se explicó en WWI-01.
\end{itemize}
\begin{good}[Resultado realmente observado]
Revisión de necesidad \#5: 6 directos, 2 derivables, 4 advertencias y 0 no disponibles. Una respuesta del asesor. Al aceptar los cuatro límites se habilitó generar, pero \textbf{no se pulsó el botón ni se ejecutó otro ETL}.
\end{good}
\copilot{Resume el importe del período y aclara qué evento y qué fecha representan los datos. ¿Este expediente demuestra que se cobró ese importe?}
\begin{note}
Es un buen primer ejercicio porque limita la complejidad. La referencia numérica del expediente \#15 no debe presentarse como resultado de esta previabilidad: pertenecen a etapas e identificadores diferentes.
\end{note}
\evidence{NEED-18 y capturas 36/37. Los números de requisitos no son un conteo de indicadores materializados. La pregunta al copiloto es una práctica posterior propuesta, condicionada a generar y publicar un datamart.}

\casepage{WWI-03}{Intermedio}{Productos facturados e importe por unidad}{Nuevo texto apoyado en capacidades conciliadas del expediente \#15.}\label{wwi03}
\begin{goal}
Comparar mensualmente el importe registrado y las unidades de las líneas de factura por producto, mostrando nombres descriptivos comprobados y conservando los identificadores de factura y línea. Ordenar los productos por la suma de sus importes y calcular el importe por unidad como importe total dividido entre unidades totales. No confundir ese cociente con el promedio simple del precio unitario.
\end{goal}
\params{WideWorldImporters}{T: Evolución de ventas en el tiempo; P: Productos con mayor desempeño.}{Mensual}{Una línea de factura.}{Fecha y producto; atributos sólo si se verifican.}
\sub{Qué revisar en la propuesta}
\begin{itemize}
\item Suma de importe registrado y cantidad. En el expediente observado se usaron \code{ExtendedPrice} y \code{Quantity}.
\item Producto por relación comprobada con \code{Warehouse.StockItems}, con identificador y descripción conservados.
\item Ranking por suma agrupada; cociente de totales para el importe por unidad y control de denominador cero.
\end{itemize}
\begin{good}
En \#15 se conciliaron 227 productos y 8.950.628 unidades. El importe global por unidad fue aproximadamente 22,126206 en moneda de origen. No es el promedio simple de \code{UnitPrice}, que fue aproximadamente 45,591688.
\end{good}
\copilot{¿Qué cinco productos tienen mayor importe facturado? Para el primero explica su cantidad, importe por unidad y denominador de participación. No uses el máximo de una sola línea para ordenar productos.}
\begin{note}
Ambos promedios pueden ser numéricamente correctos y significar cosas distintas. Además, la composición del importe puede incluir impuesto mientras el precio unitario no lo incluye. No espere igualdad sólo por llevar la palabra «precio» o «unidad».
\end{note}
\evidence{Contrato y métricas conciliadas de \#114/\#15. La capacidad está documentada; el texto nuevo y la consulta propuesta no se ejecutaron nuevamente para esta guía.}

\casepage{WWI-04}{Intermedio}{Clientes e importe promedio por factura}{Nuevo caso de práctica con cálculo documental previamente probado.}\label{wwi04}
\begin{goal}
Comparar mensualmente la facturación registrada por cliente usando su nombre descriptivo mediante relaciones verificadas. Sumar el importe de las líneas de factura, contar facturas distintas y calcular el importe promedio por factura como importe total dividido entre facturas distintas. Conservar factura y línea; no interpretar el promedio de una línea como promedio de una factura ni considerar una factura como un cobro.
\end{goal}
\params{WideWorldImporters}{T: Evolución de ventas en el tiempo; C: Desempeño por cliente.}{Mensual}{Una línea de factura relacionada con el cliente del documento.}{Fecha y cliente comercial; distinguir cuenta de facturación si aparece.}
\sub{Qué revisar en la propuesta}
\begin{itemize}
\item Importe con Suma; \code{InvoiceID} con Conteo distinto; cociente de ambos con idéntico filtro.
\item Población de clientes vinculada a factura. \code{BillToCustomerID} puede representar una cuenta distinta: no intercambiarla silenciosamente con el cliente.
\item Nombre descriptivo, cobertura y linaje sin multiplicar líneas o documentos.
\end{itemize}
\begin{good}
El expediente \#15 tiene 663 clientes descriptivos. El ticket global conciliado es 2.808,73; con 2015, 2.790,57. Son controles del expediente histórico, no valores que deba forzar en cualquier propuesta nueva.
\end{good}
\copilot{Dentro de este período, ¿qué clientes concentran más facturación y cuál es su promedio por factura? Explica qué población de cliente y documentos estás utilizando.}
\begin{note}
Si filtra líneas por producto, el promedio describe el importe seleccionado de las facturas que contienen ese producto, no necesariamente el total completo de cada factura. No sume identificadores ni use cantidad de unidades como número de clientes.
\end{note}
\evidence{QA-15/16 y expediente \#15. La dimensión y el promedio están comprobados; esta necesidad por cliente sigue siendo un ensayo propuesto hasta que se registre su propia ejecución.}

\casepage{WWI-05}{Comparación de eventos}{Pedidos registrados, no facturas emitidas}{Texto nuevo simplificado; alcance histórico probado en \#96/\#13.}\label{wwi05}
\begin{goal}
Analizar mensualmente los pedidos de venta registrados a nivel de línea. Comparar cantidades pedidas e importe calculado como precio unitario por cantidad, por producto. Contar pedidos distintos y conservar identificadores de pedido y línea. No presentar estos resultados como facturas emitidas ni como cobros.
\end{goal}
\params{WideWorldImporters}{T: Evolución de ventas en el tiempo; P: Productos con mayor desempeño.}{Mensual}{Una línea de pedido registrado, no una línea de factura.}{Fecha de pedido y producto.}
\sub{Qué revisar en la propuesta}
\begin{itemize}
\item Referencias observadas: \code{Sales.OrderLines}, \code{Sales.Orders} y \code{OrderDate}.
\item Suma de cantidades; suma del resultado precio por cantidad; pedidos distintos por \code{OrderID}.
\item Etiquetas del panel que declaren pedidos. No añadir impuestos, descuentos o estados por deducción no comprobada.
\end{itemize}
\begin{good}
El expediente histórico \#13 concilió 231.412 líneas, 73.595 pedidos, 9.310.904 unidades y 177.634.276,40 de importe según la receta de pedidos. La cabecera analítica advierte su alcance.
\end{good}
\copilot{¿Este datamart representa pedidos o facturas? Explica por qué su total no puede compararse directamente con el expediente de facturación como si fuera crecimiento comercial.}
\begin{note}
La diferencia con 198.043.439,45 del expediente de facturas no es automáticamente deuda pendiente, crecimiento ni error. Difieren evento, población y composición del importe. No reste ambos como saldo de cobranza.
\end{note}
\evidence{QA-08, QA-11, QA-13 y QA-16. El historial \#13 contiene más decisiones que este texto simplificado; volver a generar esta necesidad es una práctica nueva. No use sus costos de referencia como prueba de rentabilidad histórica.}

\casepage{WWI-06}{Intermedio condicional}{Facturación y geografía del cliente}{Nuevo texto para estudiar relaciones y nivel territorial.}\label{wwi06}
\begin{goal}
Comparar mensualmente el importe registrado y las unidades de las líneas de factura por ubicación geográfica del cliente, usando únicamente relaciones declaradas y conservando factura y línea. Mostrar un nombre territorial descriptivo e indicar si corresponde a ciudad, estado o provincia, país u otro nivel. No atribuir la ubicación a un vendedor ni convertirla en territorio comercial si esa relación no existe.
\end{goal}
\params{WideWorldImporters}{T: Evolución de ventas en el tiempo; G: Desempeño por territorio; C: Desempeño por cliente.}{Mensual}{Una línea de factura; geografía alcanzada sin duplicar el hecho.}{Fecha, cliente y nivel geográfico verificado.}
\sub{Qué revisar en la propuesta}
\begin{itemize}
\item Ruta desde factura a cliente y geografía comprobada; no elegir una ciudad desconectada sólo porque tiene nombre.
\item Aclarar qué dirección se usa y si su valor representa la ubicación actual o hay evidencia histórica. No asumir historia de cambios de domicilio.
\item Importes/unidades con Suma. La definición territorial debe ser uniforme dentro del gráfico.
\end{itemize}
\begin{good}
El expediente \#15 usa territorio descriptivo \code{StateProvinceName} y concilió 53 territorios. Esto respalda una dimensión geográfica de esa ejecución, no una regla universal de que toda fuente deba analizarse por estados.
\end{good}
\copilot{¿Qué representa el territorio de este resultado y mediante qué relación se asocia a la factura? Compara los principales territorios sin atribuir causas ni responsabilidad a vendedores.}
\begin{note}
Cambiar el nombre «estado» por «hemisferio» no modifica el nivel de los datos. Una geografía más detallada no siempre es mejor: debe responder a la decisión y conservar cobertura y comparabilidad.
\end{note}
\evidence{Dimensión y cobertura registradas en \#15; esta necesidad y la consulta sobre geografía son propuestas para ensayar. La prueba negativa sobre clima se incluye más adelante para enseñar los límites de estas relaciones.}

\chapterpage{Módulo 5 · Perfil técnico}{Arquitectura: dónde ocurre cada trabajo}
Esta sección es una ruta de entrada para el desarrollador y el operador técnico. El usuario final puede pasar al módulo de demostración. Aprender la arquitectura ayuda a diagnosticar sin tocar datos innecesariamente.
\begin{tabularx}{\linewidth}{@{}p{36mm}Y@{}}
\toprule
\textbf{Componente}&\textbf{Responsabilidad y tecnología actual}\\
\midrule
Navegador&React y TypeScript. Vite en desarrollo; Nginx para el artefacto compilado. Presenta pasos, decisiones, permisos y resultados.\\
API&FastAPI/Uvicorn. Autoriza, valida contratos y coordina módulos. El navegador no debe decidir por sí solo qué operación está permitida.\\
SQL Server&Fuentes OLTP de sólo lectura. La conexión identifica cada base. Acceso mediante pyodbc y Microsoft ODBC Driver 18.\\
PostgreSQL&Seguridad, configuración, metadatos, propuestas y auditoría en el esquema de aplicación; datos analíticos versionados por ejecución.\\
Proveedor LLM&Cloud configurado u Ollama opcional. Interpreta y devuelve estructura; no recibe permiso general para ejecutar SQL o aprobar.\\
Migraciones&SQLAlchemy/Alembic y semillas controladas. Evolucionan la aplicación sin convertir un arranque en un borrado del historial.\\
\bottomrule
\end{tabularx}
\sub{Flujo que debe poder dibujar}
El navegador envía una solicitud autenticada a la API. La API obtiene metadatos del origen o del repositorio de instantáneas; si la tarea necesita IA, prepara un contexto acotado y valida la respuesta. La aprobación se registra. Más tarde, el ETL lee el origen y escribe el datamart; la analítica consulta esa ejecución.

El origen SQL Server no es el destino analítico PostgreSQL. El esquema físico de cada ejecución sigue el patrón \code{mart_ventas_e{id}}. La selección de fuente no combina automáticamente bases ni sus catálogos.
\sub{Cuatro identificadores que evitan cruces}
\textbf{Conexión} identifica origen; \textbf{instantánea} identifica estructura; \textbf{propuesta} identifica diseño aprobado; \textbf{ejecución} identifica carga y evidencia. Mantenga su relación en peticiones, caché, estado visual, consultas y exportaciones.
\begin{note}[Un fallo visual no prueba un fallo físico, ni lo descarta]
Una respuesta tardía del navegador puede mostrar un contexto anterior. La consulta analítica debe validar también la fuente en servidor. Para diagnosticar, contraste identidad de la petición, respuesta y destino físico antes de concluir que se mezclaron datos.
\end{note}

\chapterpage{Módulo 5 · Entorno y operación técnica}{Arranque seguro y diagnóstico básico}
La ruta local de este repositorio es \path{/home/ceo/Proyectos/BI}. En otro equipo sustitúyala por la ubicación real. Lea primero \path{AGENTS.md}, alcance, arquitectura, guía de réplica y última bitácora. Los comandos del proyecto usan Docker; requieren un entorno autorizado y disponible.
\sub{Primera preparación o revisión del entorno}
\begin{verbatim}
cd "/home/ceo/Proyectos/BI"
make env
\end{verbatim}
\textbf{Antes del primer arranque:} revise localmente el archivo de entorno y cambie las contraseñas de ejemplo mediante el procedimiento autorizado. No las incluya en Git, esta guía ni capturas. \code{make env} no debe tratarse como permiso para sobrescribir una configuración existente.
\begin{verbatim}
make bootstrap
make ps
make doctor
\end{verbatim}
\code{bootstrap} valida, construye e inicia servicios y comprueba disponibilidad. Migraciones y semillas preparan la aplicación de forma controlada. Una base comercial adicional debe estar restaurada y accesible; el formulario de conexiones no restaura copias de seguridad.
\begin{tabularx}{\linewidth}{@{}Yp{42mm}p{42mm}@{}}
\toprule
\textbf{Servicio}&\textbf{Desarrollo, host}&\textbf{Entrega completa, host}\\
\midrule
Interfaz&5173&8080\\
API&8000&18000\\
PostgreSQL&55432&55433\\
SQL Server&51433&51434\\
\bottomrule
\end{tabularx}
Dentro de la red Docker de desarrollo, la aplicación usa normalmente \code{sqlserver:1433}, no el puerto publicado del host. Verifique la configuración real antes de registrar una conexión.
\sub{Comprobaciones sin borrar datos}
\begin{verbatim}
docker compose ps
docker compose logs --tail=200 backend
docker compose logs --tail=200 frontend
curl http://localhost:8000/api/v1/health/live
curl http://localhost:8000/api/v1/health/ready
\end{verbatim}
Sanee registros antes de compartirlos. Para detener preservando volúmenes existe \code{docker compose down}; \textbf{no añada opciones de eliminación de volúmenes como solución rutinaria}. Eliminar persistencia puede destruir datos o impedir descifrar secretos.

\chapterpage{Módulo 5 · Ubicación del código}{Qué archivo revisar según el problema}
{\small
\begin{tabularx}{\linewidth}{@{}p{40mm}Y@{}}
\toprule
\textbf{Tema}&\textbf{Ruta desde la raíz del repositorio}\\
\midrule
Composición API&\path{backend/app/main.py}; \path{backend/app/api/v1/router.py}\\
Usuarios, permisos, auditoría&\path{backend/app/modules/security/}\\
Conexiones y secretos&\path{backend/app/modules/parameters/}\\
Transporte de proveedores&\path{backend/app/modules/parameters/providers.py}\\
Lectura de estructura&\path{backend/app/modules/metadata/connectors.py}; \path{introspection.py}; \path{schemas.py} en ese módulo.\\
Necesidades y revisión IA&\path{backend/app/modules/copilot/need_advisor.py}; \path{needs.py} en ese módulo.\\
Propuestas y validación&\path{backend/app/modules/copilot/service.py}; \path{router.py}; \path{schemas.py} en ese módulo.\\
Nombres descriptivos&\path{backend/app/modules/copilot/entity_resolution.py}\\
Recetas y ejecución ETL&\path{backend/app/modules/etl/service.py}; \path{materializer.py}; \path{router.py} en ese módulo.\\
Panel, agregados y chat&\path{backend/app/modules/analytics/service.py}; \path{router.py} en ese módulo.\\
Exportaciones&\path{backend/app/modules/reports/analytics_export.py}\\
Pantallas&\path{frontend/src/App.tsx}; \path{AnalyticsPage.tsx}; estilos asociados.\\
Migraciones y semillas&\path{backend/migrations/versions/}; \path{backend/app/db/seed.py}\\
Contenedores y verificación&\path{compose.yaml}; \path{compose.release.yaml}; \path{Makefile}; \path{.github/workflows/}\\
\bottomrule
\end{tabularx}}
\sub{Cómo investigar sin debilitar el sistema}
Reproduzca con una entrada mínima y registre el error original. Aísle interfaz, API, proveedor o fuente. Cree una prueba que falle por la causa encontrada y luego corrija. Conserve versiones aprobadas y no cambie las reglas de negocio para acomodar una salida particular de un modelo.
\begin{note}[Frases de interfaz frente a comportamiento]
En el corte revisado quedó una ayuda de ETL con redacción de reemplazo del esquema, aunque la implementación usa destinos versionados por ejecución. También hay una discrepancia de nombre de permiso en el acta y una diferencia de PDFs entre verificación local y CI. Son aspectos de mantenimiento documental, no motivos para borrar datos ni ignorar permisos; confirme el contrato y el código vigente.
\end{note}

\chapterpage{Módulo 5 · Contratos y pruebas}{Qué no debe romper una modificación}
\begin{enumerate}
\item \textbf{Metadatos canónicos:} procedencia, tablas, tipos, claves y todas las parejas de una FK compuesta. No incluya credenciales ni muestras de filas en ese documento.
\item \textbf{Huella estable:} un orden de serialización reproducible permite verificar el artefacto. Su coincidencia no demuestra que un LLM sea determinista ni que una interpretación financiera sea verdadera.
\item \textbf{Revisión de necesidad:} objetivo, preguntas, periodicidad, instantánea y destino consentido forman el contexto. Una respuesta tardía o de otro contexto no debe habilitar generación.
\item \textbf{Propuesta controlada:} referencias existentes, tipos compatibles, rutas sin multiplicar el hecho, cobertura de lo solicitado y aprobación explícita.
\item \textbf{Ejecución y analítica:} selección de KPI, huellas, origen, destino físico y conciliación; validar que la ejecución pertenece a la fuente y puede consultarse.
\item \textbf{Historia:} consultar no debe reescribir o retirar decisiones aprobadas. Una nueva versión conserva su procedencia.
\end{enumerate}
\sub{Comandos existentes de comprobación}
\begin{verbatim}
make compose-check
make workflow-test
make test
make lint
make format-check
make verify
\end{verbatim}
\code{make test} ejecuta etapas Docker de pruebas y controles: backend con Ruff, formato, Mypy y pytest; frontend con ESLint, pruebas y compilación TypeScript/Vite. \code{make verify} añade configuración, política de ramas y documentos. No equivale a llamar de nuevo a todos los LLM ni ejecutar todos los ETL reales.
\sub{Pruebas que debe cubrir un cambio de BI + IA}
JSON inválido o incompleto; cuota y tiempo agotado; referencia inventada; tipos incompatibles; relaciones desconectadas o compuestas; precio usado como costo; tasa como dinero; denominador incorrecto; cambio de fuente con respuestas tardías; permiso insuficiente; conciliación; exportación con filtros.
\begin{good}[Resultado previo, no ejecución nueva de este manual]
La ampliación auditada cerró con 231 pruebas backend y 34 frontend, además de controles de compilación. Esta guía no reejecutó ese bloque ni usa su redacción como prueba de una nueva versión del software. Registre siempre el corte de código al repetirlo.
\end{good}

\chapterpage{Módulo 5 · Extensibilidad}{Añadir fuentes o modelos de forma general}
\sub{Otra base SQL Server no es otro motor}
Con una base SQL Server accesible, el flujo está diseñado para registrar una conexión, capturar sus metadatos, revisar su catálogo y recorrer el proceso. Aun así necesita una prueba integral: esquema leído no equivale a todos los requisitos BI soportados.
\sub{Agregar un motor distinto exige más trabajo}
\begin{enumerate}
\item Especificar conexión, credenciales, sólo lectura y driver en el entorno.
\item Normalizar el catálogo del motor: tipos, PK, FK, claves compuestas y huellas; registrar el lector de metadatos.
\item Adaptar campos y tipos permitidos en API/interfaz.
\item Adaptar la extracción y el dialecto ETL, además de comprobaciones locales de nombres, moneda y conciliación. El materializador actual usa elementos específicos de SQL Server.
\item Probar hasta panel y copiloto con un esquema distinto; incluir nombres en español/inglés, ambigüedad, nulos y relaciones que duplicarían hechos.
\end{enumerate}
Agregar sólo un lector de tablas no hace que MySQL u otro motor esté soportado de extremo a extremo. Evite condiciones por nombre de base como «si es AdventureWorks»: use capacidades y evidencia del contrato.
\sub{Agregar otro proveedor LLM}
\begin{enumerate}
\item Especificar autenticación, URL permitida, capacidades, razonamiento y límites.
\item Implementar transporte y extracción de respuesta en el adaptador, con prueba de conexión y errores seguros.
\item Mantener validación local de formato y negocio. Reintentos acotados, sin cambio silencioso del destino autorizado.
\item Añadir pruebas de éxito, cuota, timeout, JSON inválido, credencial ausente y modelo no disponible.
\item Repetir casos reales autorizados registrando modelo y resultados; no cambiar fórmulas financieras por proveedor.
\end{enumerate}
\begin{note}[Principio de diseño]
La adaptación del protocolo es legítima; relajar la semántica para que un modelo «pase» no lo es. Una referencia desconocida sigue siendo desconocida aunque el proveedor afirme con mucha seguridad que es correcta.
\end{note}

\chapterpage{Módulo 5 · Mantenimiento responsable}{Cambios, documentación y recuperación}
\sub{Desarrollo guiado por especificaciones}
Antes de cambiar comportamiento, actualice alcance y criterios en la especificación correspondiente. Si altera una decisión de arquitectura duradera, documente la decisión. Añada pruebas, implemente, verifique y actualice manuales, bitácora y casos de QA. No convierta un incidente aislado en una excepción fija para una base de ejemplo.

El flujo de integración es rama de trabajo, PR a \code{develop} y promoción revisada hacia \code{main}. Este manual no crea PR ni modifica ramas. Una promoción debe identificar qué se verificó y qué queda pendiente.
\sub{Generación documental disponible en el proyecto}
\begin{verbatim}
make technical-manual
make qa-audit-report
make thesis-chapter3
make docs
\end{verbatim}
Revise visualmente los PDF; compilar no prueba que tablas y capturas sean legibles. En el corte revisado, \code{make docs} cubre once PDF, mientras el flujo CI enumeraba diez y no incluía aún el PDF de auditoría del 3 de octubre. No afirme paridad total sin corregir y comprobar esa diferencia.
\sub{Entornos que no deben confundirse}
\textbf{Desarrollo:} \code{make up}. \textbf{Vista compilada local:} \code{make delivery-preview-up}, que comparte API y datos de desarrollo. \textbf{Entrega completa:} \code{make release-up}, con configuración y volúmenes propios. No levante dos modalidades que disputen el mismo puerto sin revisar su configuración.

\textbf{Ollama opcional:} existen \code{make ollama-up}, \code{make ollama-pull}, \code{make ollama-status} y \code{make ollama-down}. Descargar un modelo consume red/disco y no lo convierte automáticamente en el proveedor activo. Su capacidad debe comprobarse como la de cualquier otro modelo.
\sub{Secretos, respaldo y recuperación}
Las credenciales LLM se registran en la aplicación, no en código ni capturas. La persistencia cifrada depende de su clave de descifrado local: perder esa raíz puede impedir recuperar secretos. El responsable debe conservar respaldos de datos y material de recuperación según política, comprobar restauraciones y no publicar esos respaldos como evidencia de tesis.
\begin{good}[Secuencia frente a un incidente técnico]
Identificar → preservar evidencia y datos → reproducir en alcance controlado → corregir con prueba → verificar regresión → actualizar documentación → comunicar límites. Reiniciar servicios puede recuperar disponibilidad, pero no demuestra que la causa haya sido corregida.
\end{good}

\chapterpage{Módulo 6 · Preparar la exposición}{Guion de demostración de 25 minutos}
Prepare una ruta principal con expedientes existentes y una alternativa de práctica en vivo. Así podrá explicar el sistema incluso si se agota la cuota del proveedor. No presente una consulta histórica como generación nueva.
\begin{tabularx}{\linewidth}{@{}p{18mm}YY@{}}
\toprule
\textbf{Minutos}&\textbf{Acción visible}&\textbf{Explicación sugerida}\\
\midrule
0--3&Presentar problema y conceptos.&«Convertimos una necesidad comercial en un diseño verificable. La IA asiste; la plataforma comprueba; una persona aprueba».\\
3--6&Explorador AW y WWI.&«Los metadatos describen fuentes diferentes. Una línea de pedido no es una línea de factura».\\
6--8&Catálogos por fuente.&«Las preguntas orientan, pero no inventan columnas. AW conserva sus personalizaciones; WWI tiene su propio catálogo».\\
8--12&Necesidad, consentimiento y viabilidad.&«Los derivables son candidatos. La estructura y los límites se revisan antes de generar». Usar capturas si no se hará una llamada.\\
12--16&Abrir \#98 y \#114.&«Estas propuestas conservan grano y documentos. El promedio por pedido/factura usa un conteo distinto».\\
16--19&Abrir expedientes \#14 y \#15.&«Aprobar no fue cargar. Aquí están la ejecución y la conciliación de origen y destino».\\
19--23&Panel, filtros y gráfico interactivo.&«El dato cambia con el filtro y conserva procedencia. El total de WWI no se supone USD».\\
23--25&Copiloto, exportación y límites.&«Contrasto indicador, evento, filtro y denominador. No afirmo causalidad ni compatibilidad no probada».\\
\bottomrule
\end{tabularx}
\sub{Lista de preparación del expositor}
\begin{itemize}
\item Probar acceso y permisos; confirmar ambas fuentes sin cambiar sus catálogos.
\item Localizar las propuestas y los expedientes antes de exponer. Si sus datos físicos ya no están, no prometer navegación en vivo del panel.
\item Preparar capturas y exportaciones con fecha, fuente y filtros. No mostrar claves.
\item Comprobar cuota sólo si se planean nuevas llamadas autorizadas. No prometer una latencia exacta.
\item Si se genera un caso nuevo, registrar su propio número. No reemplazar los identificadores de la evidencia histórica.
\end{itemize}
\textbf{Cierre sugerido:} «Este recorrido está acreditado en dos fuentes SQL Server. Rechazar una sugerencia sin respaldo y conservar evidencia es parte de la calidad del sistema».

\chapterpage{Módulo 6 · Aprendizaje activo}{Taller guiado para los tres perfiles}
Los tiempos son orientativos. La práctica de 80 minutos supone aplicación, fuentes, cuenta y proveedor preparados. El estudio de conceptos puede realizarse antes, sin consumir IA.
\begin{tabularx}{\linewidth}{@{}p{19mm}Yp{45mm}@{}}
\toprule
\textbf{Tiempo}&\textbf{Actividad}&\textbf{Producto del alumno}\\
\midrule
0--10&Leer conceptos, iniciar sesión y reconocer la fuente.&Explicar pedido, factura, fila y KPI con sus palabras.\\
10--20&Explorar ambas bases y una relación real.&Identificar PK, FK, fecha y grano candidato.\\
20--30&Comparar catálogos y escribir objetivo pequeño.&Texto, preguntas, periodicidad y límites.\\
30--45&Autorizar y validar; leer cada resultado.&Registro de directos, derivables y decisiones pendientes.\\
45--55&Generar si está autorizado y revisar contrato.&Grano, medida, fórmula, denominador y cobertura.\\
55--65&Aprobar y ejecutar un caso controlado, o consultar histórico si no alcanza el tiempo.&Identificar honestamente ejecución nueva o consulta.\\
65--75&Filtrar, abrir gráfico, preguntar y exportar.&Comparación de cifras y captura del alcance.\\
75--80&Redactar conclusión e incidente si corresponde.&Acta con esperado, obtenido, evidencia y límite.\\
\bottomrule
\end{tabularx}
\sub{Ruta por perfil}
\textbf{Usuario final / gerencia:} estudie conceptos y criterios de ventas; consulte panel, copiloto y exportaciones. Debe poder cuestionar un promedio sin administrar credenciales.

\textbf{Operador / analista autorizado:} añada catálogo, necesidad, diseño, personalización, aprobación, ETL y conciliación. Debe reconocer cuándo detener el proceso.

\textbf{Administrador / desarrollador:} añada entorno, seguridad, arquitectura, contratos, diagnóstico y pruebas. Debe resolver causas sin debilitar los controles ni confundir un problema de interfaz con uno de datos.
\begin{note}[Condición de aprobación del taller]
No se aprueba por llegar a un gráfico. El alumno debe identificar evento, grano, fuente, fórmula y límites, y saber dónde está la conciliación. Un caso negativo bien diagnosticado también demuestra aprendizaje; no debe disfrazarse de caso positivo.
\end{note}
\sub{Ejercicio de explicación entre compañeros}
Una persona hace de gerente y pregunta «¿puedo usar esto para decidir a qué cliente dar crédito?». La otra debe explicar que ventas o facturas por sí solas no prueban capacidad de pago: faltan datos y criterios adecuados. No invente una decisión crediticia desde un ranking comercial.

\chapterpage{Módulo 6 · Valores de referencia}{Controles para repetir sin confundir poblaciones}
Estos resultados se observaron en la auditoría del 3 de octubre de 2026 o en el expediente histórico señalado. \textbf{No son resultados nuevos producidos al crear este manual}. Identificadores y cifras dependen de la instalación y su corte.
\begin{tabularx}{\linewidth}{@{}Yrrr@{}}
\toprule
\textbf{Control}&\textbf{AW pedidos}&\textbf{WWI pedidos}&\textbf{WWI facturas}\\
\midrule
Propuesta / ejecución&98 / 14&96 / 13&114 / 15\\
Filas de detalle&121.317&231.412&228.265\\
Documentos distintos&31.465&73.595&70.510\\
Unidades&No se usa aquí&9.310.904&8.950.628\\
\bottomrule
\end{tabularx}
\sub{Importes y promedios que sí puede contrastar}
\begin{tabularx}{\linewidth}{@{}Yp{53mm}@{}}
\toprule
\textbf{Contexto}&\textbf{Valor observado}\\
\midrule
AW \#14: importe de todas las líneas&109.846.381,4250\\
AW \#14: promedio por pedido&3.491,07\\
AW \#14, 2014: pedidos / promedio&11.761 / 1.705,46\\
WWI \#13: importe según receta de pedidos&177.634.276,40\\
WWI \#15: importe facturado registrado&198.043.439,45\\
WWI \#15: promedio por factura&2.808,73\\
WWI \#15, 2015: total / facturas&62.090.220,81 / 22.250\\
WWI \#15, 2015: promedio por factura&2.790,57\\
\bottomrule
\end{tabularx}
\sub{Cuatro promedios de WWI que no deben confundirse}
En el expediente \#15: importe medio por línea aproximadamente \textbf{867,603178}; precio unitario medio simple \textbf{45,591688}; importe por unidad \textbf{22,126206}; importe por factura \textbf{2.808,728399}. La etiqueta, la composición del importe y el denominador hacen que respondan preguntas distintas.
\begin{note}[No es una comparación monetaria entre empresas]
AW mostró USD en el expediente observado; WWI conserva moneda de origen sin código único demostrado. Además, pedidos y facturas difieren. No calcule diferencias, participaciones o crecimiento entre las tres columnas como si fueran una sola serie de ventas comparable.
\end{note}
\textbf{Regla de contraste:} use el valor detallado y la precisión del expediente. Una tarjeta abreviada, como 109,8 M, no permite comprobar diferencias de centavos.

\chapterpage{Módulo 6 · QA para el alumno}{Pruebas positivas, negativas y registro de evidencia}
\begin{tabularx}{\linewidth}{@{}p{17mm}YY@{}}
\toprule
\textbf{Caso}&\textbf{Qué hacer}&\textbf{Resultado que debe comprobar}\\
\midrule
P-01&Consultar AW \#14 y filtrar un año presente.&Fuente/ejecución correctas; valores y gráficos responden al mismo filtro.\\
P-02&Consultar WWI \#15 y filtrar 2015.&Facturas, no pedidos; cifras coherentes con su expediente y denominador.\\
P-03&Volver de WWI a AW.&No quedan productos, territorios, revisión ni chat de la otra fuente como resultados propios.\\
N-01&Pedir clima externo con el texto siguiente.&Límite visible; no inventa clima ni sustituye meteorología por sensores internos.\\
N-02&Tras una revisión, editar el objetivo.&La revisión anterior no autoriza silenciosamente el alcance nuevo; revisar consentimiento.\\
N-03&Revisar una propuesta cuyo promedio sea de línea.&No aprobarlo como promedio documental. Corregir o etiquetar con honestidad.\\
N-04&Comparar factura con cobro en el copiloto.&No afirmar cobro sin evidencia; declarar limitación.\\
\bottomrule
\end{tabularx}
\begin{goal}
Analizar las ventas mensuales registradas y explicar cuánto influyó el clima de cada ciudad usando temperaturas externas que no están en esta base de datos.
\end{goal}
El negativo climático fue observado en NEED-20: se detectó el dato ausente y no se aceptó alcance parcial ni se generó ETL. El ensayo también encontró un problema temporal corregido antes del positivo final. \textbf{Un nuevo intento debe documentarse aparte}; no se ejecuta para demostrar que todo texto debe pasar.
\sub{Plantilla de ficha para la tesis o el taller}
\begin{tabularx}{\linewidth}{@{}p{42mm}Y@{}}
\toprule
Identificación&ID del caso, fecha, responsable y perfil utilizado.\\
Preparación&Fuente, instantánea, proveedor/modelo, permisos y versión del programa.\\
Entradas&Objetivo exacto, preguntas, periodicidad y límites aceptados.\\
Pasos&Acciones reproducibles; indicar consulta histórica o generación nueva.\\
Esperado / obtenido&Criterio antes de probar; observación real después, sin reescribir el esperado para hacerlo pasar.\\
Evidencia&Versión/ejecución si existen, capturas, filtros y conciliación.\\
Conclusión&Aprobado, fallido, parcial o no ejecutado; incidencia y riesgo pendiente.\\
\bottomrule
\end{tabularx}

\chapterpage{Módulo 6 · Qué hacer ante problemas}{Acceso, preparación e inteligencia artificial}
\sub{Primero deténgase, identifique y conserve}
No pulse varias veces una operación con efecto real. Registre pantalla, mensaje, fuente, versión, hora y última acción. No incluya credenciales. Distinga si falló el \textbf{acceso}, la \textbf{conexión}, el \textbf{proveedor}, el \textbf{contrato} o el \textbf{cálculo}.
{\small
\begin{tabularx}{\linewidth}{@{}p{29mm}YY@{}}
\toprule
\textbf{Síntoma}&\textbf{Acción segura del usuario}&\textbf{Cuándo escalar / qué evitar}\\
\midrule
No abre la página&Verificar dirección y si está en el equipo correcto; avisar al responsable técnico.&No borrar volúmenes ni reinstalar por intuición. Revisar servicios y puertos.\\
No puede entrar&Comprobar correo; usar el acceso asignado. Si venció sesión, volver a autenticar.&No compartir claves ni usar otra cuenta para saltarse permisos.\\
No aparece una opción&Confirmar perfil y permiso con administrador.&La vista ejecutiva no otorga permisos. No modificar roles para una demostración sin autorización.\\
Fuente no lista&Revisar conexión probada, habilitación e instantánea.&El administrador revisa acceso de sólo lectura, servidor y base. No duplicar conexión por cada fallo.\\
No hay metadatos&Crear primera instantánea con permiso. Si hay una, revisar contexto y versión.&Actualizar estructura no carga ventas. No cambiar metadatos para que acepte una referencia falsa.\\
Validar deshabilitado&Objetivo entre 20 y 2.000 caracteres, pregunta, período, destino y consentimiento.&Leer mensaje; adoptar una sugerencia obliga a revisar otra vez el contexto.\\
Sugerencia bloqueada&Mantener texto propio; revisar qué referencia o capacidad falta.&No asumir que la necesidad comercial es inválida sólo porque una reformulación falló.\\
Cuota / límite temporal&Conservar objetivo; esperar el intervalo indicado o usar históricos.&No reintentar sin límite ni cambiar proveedor sin coordinar la configuración activa.\\
JSON inválido o truncado&No usar la respuesta descartada. Reintentar de forma acotada cuando la aplicación lo permita.&Desarrollador revisa categoría del error y contrato. No copiar JSON parcial para aprobarlo manualmente.\\
Necesidad no disponible&Reducir alcance con decisión explícita o aportar una fuente autorizada.&No inventar datos. Si lo ausente es central para la decisión, no continúe como si estuviera resuelto.\\
\bottomrule
\end{tabularx}}
\begin{note}[Durante una exposición]
Diga qué falló, muestre que se conservó el trabajo y pase al expediente de referencia identificado como histórico. Esa recuperación es demostrable; ocultar el fallo y afirmar que la generación fue exitosa no lo es.
\end{note}

\chapterpage{Módulo 6 · Qué hacer ante problemas}{Diseño, ETL, panel y coherencia comercial}
{\small
\begin{tabularx}{\linewidth}{@{}p{29mm}YY@{}}
\toprule
\textbf{Síntoma}&\textbf{Qué comprobar y hacer}&\textbf{Qué no debe hacer}\\
\midrule
Error bloqueante en propuesta&Leer referencias, grano, tipos y cobertura. Personalizar una opción válida o generar una nueva versión corregida.&No aprobar ni eliminar la regla para desbloquear.\\
Histórico con advertencias&Distinguir aprobación histórica, compatibilidad actual y datos disponibles. Verificar evidencia y registrar diferencias.&No retirar aprobación ni descartar versiones sólo por consultarlas.\\
No puede crear otra necesidad&Usar Crear nueva propuesta desde el histórico.&No cerrar sesión como método normal ni borrar el resultado anterior.\\
KPI Sin valor&Revisar estado, fórmula, medida base, denominador y filtro.&No sustituir nulo por cero ni afirmar «no hubo ventas» sin respaldo.\\
Promedio incoherente&Identificar si es por línea, documento o unidad. Comparar totales del mismo alcance.&No cambiar sólo la etiqueta para justificar otra fórmula.\\
ETL fallido&Abrir expediente, leer etapa y error; técnico revisa registros y conectividad.&No ejecutar repetidamente ni eliminar datamarts anteriores.\\
Diferencia en conciliación&Conservar resultados, revisar población, uniones, precisión y fórmula antes de publicar.&No normalizar la diferencia como aceptable sin criterio documentado.\\
Etiqueta pendiente&Revisar mapeos y usar publicación o reintento de interpretación sin repetir ETL.&No traducir identificadores o nombres personales ni cargar de nuevo por un texto.\\
Moneda desconocida&Usar Comprobar divisa sin repetir ETL si aparece; mantener moneda de origen si no se prueba.&No asumir USD ni convertir importes sin evidencia.\\
Panel ausente&Comprobar fuente, permiso, ejecución conciliada/publicada y tablas físicas.&Un historial guardado no garantiza datos físicos consultables.\\
Gráfico y chat distintos&Comparar evento, indicador, filtros, Top N, denominador y redondeo.&No asumir cruce de fuentes antes de revisar el alcance interpretado; si hay cruce real, detener análisis.\\
Fuente cruzada visible&Registrar pasos/captura y detener conclusiones; escalar a QA/desarrollo.&No ocultar el problema cambiando etiquetas ni reconstruir a ciegas ambas fuentes.\\
\bottomrule
\end{tabularx}}
\sub{Paquete mínimo para solicitar ayuda}
Pantalla y acción; fuente/instantánea; propuesta/ejecución; objetivo sin secretos; indicador y filtros; esperado y obtenido; fecha; captura legible. El técnico puede añadir categoría de error y registros saneados. Un caso reproducible vale más que «no funciona».

\chapterpage{Módulo 6 · Evaluación de aprendizaje}{Preguntas para comprobar que comprendió BI + IA}
Intente responder antes de leer la siguiente página. No hace falta memorizar código.
\begin{enumerate}
\item ¿En qué se diferencia una fuente OLTP de un datamart?
\item ¿Qué es el grano? Si una factura tiene tres líneas, ¿cuántas facturas cuenta?
\item ¿Qué diferencia existe entre dimensión Producto y atributo Categoría?
\item ¿Qué demuestra una FK y qué no demuestra sobre los importes?
\item ¿Por qué Mensual no significa perder el detalle ni programar ETL mensual?
\item ¿Pedido, factura y cobro pueden usarse como sinónimos? Dé un contraejemplo.
\item ¿Cuándo no debe llamar venta neta al importe de una columna?
\item ¿Cuál es la diferencia entre promedio de línea, por documento y por unidad?
\item ¿Por qué el producto con la línea más cara puede no ser el producto líder?
\item ¿Qué diferencia existe entre costo estándar, último costo y costo histórico?
\item ¿Qué hace la IA al sugerir una necesidad y qué sigue decidiendo la persona?
\item ¿Qué significa derivable? ¿Ya está conciliado ese cálculo?
\item ¿Aprobar una propuesta ejecuta ETL? ¿Verificar evidencia lo ejecuta?
\item ¿Los Hallazgos explicables del panel necesitan una llamada al LLM?
\item Si una respuesta del copiloto usa otro año, ¿debe coincidir con el panel actual?
\item ¿Qué significa que las filas concilien? ¿Certifica toda la interpretación comercial?
\item ¿Qué tres datos de procedencia debe mirar antes de leer el tablero?
\item ¿Por qué «moneda de origen» puede ser más correcto que mostrar USD?
\item ¿Cambiar Vista ejecutiva a Vista analítica demuestra que otro rol tiene acceso?
\item ¿Qué debe registrar cuando una prueba no puede completarse por cuota?
\end{enumerate}
\begin{good}[Criterio de dominio]
Puede usar la herramienta con criterio cuando responde estas preguntas con ejemplos, identifica la evidencia de una cifra y sabe detener una aprobación. Saber navegar sin comprender el evento y el denominador todavía no es dominar BI.
\end{good}

\chapterpage{Módulo 6 · Respuestas de estudio}{Respuestas comentadas y errores que evitan}
{\small
\begin{enumerate}
\item La fuente registra operaciones; el datamart organiza un alcance analítico. Aquí SQL Server es fuente y PostgreSQL almacena la materialización.
\item El grano define una fila. Tres líneas de una factura son tres hechos de detalle, pero un documento distinto. Evita inflar conteos.
\item Producto es el conjunto descriptivo; Categoría es uno de sus atributos. No es necesario crear una dimensión independiente por cada campo.
\item Una FK acredita una relación declarada entre claves. No acredita pago, composición fiscal ni calidad económica de los valores.
\item Mensual agrupa por año y mes. La fila puede seguir siendo una línea; la agenda de cargas es otro asunto.
\item No. Un pedido puede no haberse facturado; una factura puede estar pendiente de cobro. Cada evento requiere fuente y fecha propias.
\item Cuando no se conoce si incluye descuentos, impuestos, ajustes o qué población representa. Use importe registrado y documente límites.
\item Línea: importe / filas; documento: importe / documentos distintos; unidad: importe / unidades. El filtro y la base del importe deben coincidir.
\item Una línea máxima es una observación individual; el liderazgo por producto depende de sumar sus líneas. Confundirlos cambia la decisión.
\item Estándar es una referencia; último costo describe el valor disponible más reciente; histórico corresponde al momento/evento si existe evidencia. No son intercambiables.
\item La IA propone una redacción según contexto. La persona elige objetivo, acepta límites y aprueba después; la plataforma comprueba la salida.
\item Existe una ruta o receta candidata. Debe implementarse en el contrato y luego ejecutarse y conciliarse. No equivale a dato ya calculado.
\item No en ambos casos. Aprobar habilita una propuesta elegible; verificar comprueba estructura. La carga requiere la confirmación específica de materialización.
\item No: esos resúmenes se calculan por reglas. La conversación del copiloto sí puede interpretar y redactar mediante IA.
\item Sólo si usa el mismo alcance. Lea consulta interpretada y filtros; una consulta permitida no necesariamente cambia el tablero.
\item En estos ensayos acredita igualdad del detalle cargado; se revisan también medidas/KPI. No certifica por sí sola historia de costos, cobro o interpretación fiscal.
\item Fuente, propuesta y ejecución. Luego revise indicador, unidad, filtros y fecha de la carga.
\item Porque puede no existir evidencia de una moneda ISO única. Inventar un símbolo añade una certeza falsa.
\item No. Son modos de presentación. Debe probarse una cuenta autenticada con sus permisos reales.
\item Entradas, proveedor/modelo, hora, error, estado alcanzado y alternativa utilizada. Marque parcial/no ejecutado; no lo cuente como éxito de generación.
\end{enumerate}}
\sub{Ejercicio final de síntesis}
Explique en dos minutos una cifra del panel con esta estructura: \emph{«Proviene de esta fuente y ejecución; mide este evento con este grano; usa esta fórmula y estos filtros; se concilió así; y no permite concluir esto otro»}. Si puede hacerlo, está conectando operación, BI y juicio comercial.

\chapterpage{Referencia y trazabilidad documental}{Fuentes, vigencia y uso académico}
\sub{Qué se utilizó para construir esta guía}
\begin{enumerate}[label=R\arabic*.]
\item \textbf{Auditoría operativa BI + IA, 3 de octubre de 2026.} Archivo del proyecto \path{docs/testing/auditoria-operativa-2026-10-03.md}. Casos QA-01 a QA-16 y ampliación NEED-01 a NEED-21. Sustenta resultados, incidentes, alcance y limitaciones.
\item \textbf{Manual técnico del prototipo.} \path{docs/manual-tecnico/manual-tecnico.tex}. Secciones de operación, asistencia, materialización, analítica, configuración y recuperación. Se contrastaron sus instrucciones con el código actual.
\item \textbf{Arquitectura, alcance y réplica.} \path{docs/architecture.md}, \path{docs/scope.md}, \path{docs/replication-guide.md} y \path{AGENTS.md}. Sustentan la separación de responsabilidades y el trabajo de desarrollo.
\item \textbf{Interfaz vigente.} \path{frontend/src/App.tsx} y \path{frontend/src/AnalyticsPage.tsx}. Sustentan nombres de botones, estados, rutas, consentimiento y filtros; algunas capturas históricas muestran una disposición anterior.
\item \textbf{Contratos y servicios.} Módulos backend de metadatos, parámetros/proveedores, copiloto, ETL, analítica y reportes; Makefile y flujos CI/CD. Se revisaron para no atribuir al LLM tareas determinísticas ni prometer motores no implementados.
\item \textbf{Registros persistidos consultados en sólo lectura.} Catálogos, propuestas \#96/\#98/\#114 y expedientes \#13/\#14/\#15 del entorno de demostración. Sustentan la diferenciación de eventos y los controles de agregados.
\end{enumerate}
\sub{Procedencia de las capturas}
Las doce imágenes incluidas provienen de \path{docs/testing/evidencias/qa-20261003/}. Sus nombres originales se conservan en la carpeta \path{figuras/} del paquete LaTeX. Se reutilizan como evidencia real de aquella auditoría; no son interfaces recreadas ni pruebas nuevas realizadas durante la redacción.

Las capturas cubren explorador, catálogo, consentimiento, viabilidad, interpretación IA, aprobación, etapas de ETL, conciliación, panel, conversación y distinción de pedidos. La guía explica qué demuestra cada imagen y qué no.
\sub{Qué sigue pendiente fuera de este documento}
Una ejecución nueva de cada texto propuesto; pruebas de cuentas autenticadas de los roles nuevos; validación integral de otros motores/modelos/esquemas; y cualquier afirmación histórica o financiera que requiera datos adicionales. Esta guía no transforma esos pendientes en resultados aprobados.
\begin{good}[Cómo citarla en su trabajo]
Puede identificarla como \emph{Guía de formación desde cero: BI + IA aplicada a ventas}, versión 1.0, 3 de octubre de 2026, material de apoyo del prototipo. Para afirmar un resultado experimental, cite además el caso de auditoría y su evidencia; no use una ficha propuesta como si fuera una observación obtenida.
\end{good}


\end{document}
